

\documentclass[fleqn,leqno,11pt]{article}

%%%%%%%%%%%%%%%%%%%%%%%%%%%%%%%%%%%%%%%%%%%%%%%%%%%%%%%%%%%%%%%%%%%%%%%%%%%%%%%%%%%%%%%%%%%%%%%%%%%%%%%%%%%%%%%%%%%%%%%%%%%%%%%%%%%%%%%%%%%%%%%%%%%%%%%%%%%%%%%%%%%%%%%%%%%%%%%%%%%%%%%%%%%%%%%%%%%%%%%%%%%%%%%%%%%%%%%%%%%%%%%%%%%%%%%%%%%%%%%%%%%%%%%%%%%%

\usepackage{amssymb}

\usepackage{econ}



%TCIDATA{OutputFilter=LATEX.DLL}

%TCIDATA{Version=5.50.0.2953}

%TCIDATA{<META NAME="SaveForMode" CONTENT="1">}

%TCIDATA{BibliographyScheme=Manual}

%TCIDATA{Created=Tue Feb 10 10:33:10 1998}

%TCIDATA{LastRevised=Monday, December 18, 2006 11:28:06}

%TCIDATA{<META NAME="GraphicsSave" CONTENT="32">}

%TCIDATA{Language=American English}

%TCIDATA{CSTFile=Econ.cst}



\input{tcilatex}

\renewenvironment{description}

                {\list{}{\listparindent -\parindent

                  \let\makelabel\descriptionlabel

                  \itemsep       8pt

                  \parsep        0pt

                        \itemindent    \listparindent

                        \leftmargin    \parindent}}

                {\endlist}

\input{lefttag.sty}

\setstretch{1.48}

\setlength{\topmargin}{.3in}

\setlength{\oddsidemargin}{.75in}

\setlength{\evensidemargin}{.75in}

\setlength{\textwidth}{6.9in}

\setlength{\textheight}{9.05in}

\setlength{\footskip}{0.4in}

\setlength{\parindent}{0.4in}

\setcounter{page}{1}

\font\smcaps=cmcsc10



\begin{document}



\title{BUSINESS CYCLE ACCOUNTING}

\author{V. V. Chari, Patrick J. Kehoe, and Ellen R. McGrattan\thanks{%

We thank the co-editor and three referees for useful comments. We also thank

Kathy Rolfe for excellent editorial assistance and the National Science

Foundation for financial support. The views expressed herein are those of

the authors and not necessarily those of the Federal Reserve Bank of

Minneapolis or the Federal Reserve System.}}

\maketitle



\begin{abstract}

We propose a simple method to help researchers develop quantitative models

of economic fluctuations. The method rests on the insight that many models

are equivalent to a prototype growth model with time-varying wedges which

resemble productivity, labor and investment taxes, and government

consumption. Wedges corresponding to these variables---\textit{efficiency}, 

\textit{labor}, \textit{investment,} and \textit{government consumption

wedges}---are measured and then fed back into the model in order to assess

the fraction of various fluctuations they account for. Applying this method

to U.S. data for the Great Depression and the 1982 recession reveals that

the efficiency and labor wedges together account for essentially all of the

fluctuations; the investment wedge plays a decidedly tertiary role, and the

government consumption wedge, none. Analyses of the entire postwar period

and alternative model specifications support these results. Models with

frictions manifested primarily as investment wedges are thus not promising

for the study of business cycles.\newline

\end{abstract}



\noindent In building detailed, quantitative models of economic

fluctuations, researchers face hard choices about where to introduce

frictions into their models in order to allow the models to generate

business cycle fluctuations similar to those in the data. Here we propose a

simple method to guide these choices, and we demonstrate how to use it.



Our method has two components: an equivalence result and an accounting

procedure. The \textit{equivalence result} is that a large class of models,

including models with various types of frictions, are equivalent to a

prototype model with various types of time-varying wedges that distort the

equilibrium decisions of agents operating in otherwise competitive markets.

At face value, these wedges look like time-varying productivity, labor

income taxes, investment taxes, and government consumption. We thus label

the wedges \textit{efficiency wedges}, \textit{labor wedges}, \textit{%

investment wedges}, and \textit{government consumption wedges}.



The \textit{accounting procedure} also has two components. It begins by

measuring the wedges, using data together with the equilibrium conditions of

a prototype model. The measured wedge values are then fed back into the

prototype model, one at a time and in combinations, in order to assess how

much of the observed movements of output, labor, and investment can be

attributed to each wedge, separately and in combinations. By construction,

all four wedges account for all of these observed movements. This accounting

procedure leads us to label our method \textit{business cycle accounting}.



To demonstrate how the accounting procedure works, we apply it to two actual

U.S. business cycle episodes: the most extreme in U.S. history, the Great

Depression (1929--39), and a downturn less severe and more like those seen

since World War II, the 1982 recession. For the Great Depression period, we

find that, in combination, the efficiency and labor wedges produce declines

in output, labor, and investment from 1929 to 1933 only slightly more severe

than in the data. These two wedges also account fairly well for the behavior

of those variables in the recovery. Over the entire Depression period,

however, the investment wedge actually drives output the wrong way, leading

to an increase in output during much of the 1930s. Thus, the investment

wedge cannot account for either the long, deep downturn or the subsequent

slow recovery. Our analysis of the more typical 1982 U.S. recession produces

essentially the same results for the efficiency and labor wedges in

combination. Here the investment wedge plays essentially no role. In both

episodes, the government consumption wedge plays virtually no role.



We extend our analysis to the entire postwar period by developing some

summary statistics for 1959--2004. The statistics we focus on are the output

fluctuations induced by each wedge alone and the correlations between those

fluctuations and those actually in the data. Our findings from these

statistics suggest that over the entire postwar period the investment wedge

plays a somewhat larger role in business cycle fluctuations than in the 1982

recession, but its role is substantially smaller than that of either the

labor or efficiency wedges.



We begin our demonstration of our proposed method by establishing

equivalence results that link the four wedges to detailed models. We start

with detailed model economies in which technologies and preferences are

similar to those in a benchmark prototype economy and show that frictions in

the detailed economies manifest themselves as wedges in the prototype

economy. We show that an economy in which the technology is constant but

input-financing frictions vary over time is equivalent to a growth model

with efficiency wedges. We show that an economy with sticky wages and

monetary shocks, like that of Bordo, Erceg, and Evans (2000), is equivalent

to a growth model with labor wedges. In the appendix, we show that an

economy with the type of credit market frictions considered by those of

Bernanke, Gertler, and Gilchrist (1999) is equivalent to a growth model with

investment wedges. Also in the appendix, we show that an open economy model

with fluctuating borrowing and lending is equivalent to a prototype

(closed-economy) model with government consumption wedges. In the working

paper version of this paper (Chari, Kehoe, and McGrattan (2004)), we also

show that an economy with the type of credit market frictions considered by

Carlstrom and Fuerst (1997) is equivalent to a growth model with investment

wedges and that an economy with unions and antitrust policy shocks, like

that of Cole and Ohanian (2004), is equivalent to a growth model with labor

wedges.



Similar equivalence results can be established when technology and

preferences in detailed economies are very different from those in the

prototype economy. In such situations, the prototype economy can have wedges

even if the detailed economies have no frictions. We show how wedges in the

benchmark prototype economy can be decomposed into a part due to frictions

and a part due to differences in technology and preferences by constructing

alternative prototype economies which have technologies and preferences

similar to those in the detailed economy.



Our quantitative findings suggest that financial frictions which manifest

themselves primarily as investment wedges do not play a primary role in the

Great Depression or postwar recessions. Such financial frictions play a

prominent role in the models of Bernanke and Gertler (1989), Carlstrom and

Fuerst (1997), Kiyotaki and Moore (1997), and Bernanke, Gertler, and

Gilchrist (1999). More promising, our findings suggest, are models in which

the underlying frictions manifest themselves as efficiency and labor wedges.

One such model is the input-financing friction model described here in which

financial frictions manifest themselves primarily as efficiency wedges. This

model is consistent with the views of Bernanke (1983) on the importance of

financial frictions. Also promising are sticky wage models with monetary

shocks, such as that of Bordo, Erceg, and Evans (2000), and models with

monopoly power, such as that of Cole and Ohanian (2004) in which the

underlying frictions manifest themselves primarily as labor wedges. In

general, this application of our method suggests that successful future work

will likely include mechanisms in which efficiency and labor wedges have a

primary role and the investment wedge has, at best, a tertiary role. We view

this finding as our key substantive contribution.



In our quantitative work, we also analyze some detailed economies with quite

different technology and preferences than those in our benchmark prototype

economy. These include variable instead of fixed capital utilization,

different labor supply elasticities, and costs of adjusting investment. For

these alternative detailed economies, we decompose the benchmark prototype

wedges into their two sources, frictions and specification differences, by

constructing alternative prototype economies that are equivalent to the

detailed economies and so can measure the part of the wedges due to

frictions. We find that with regard to the investment wedge's role in the

business cycle, frictions driving that wedge are unchanged by different

labor supply elasticities and worsened by variable capital

utilization---with the latter specification, for example, the investment

wedge boosts output even more during the Great Depression than it did in the

benchmark economy. With investment adjustment costs, the frictions driving

investment wedges do at least depress output during the downturns, but only

modestly. Altogether, these analyses reinforce our conclusion that the

investment wedge plays a decidedly tertiary role in business cycle

fluctuations.



Our business cycle accounting method is intended to shed light on promising

classes of mechanisms through which primitive shocks lead to economic

fluctuations. It is not intended to identify the primitive sources of

shocks. Many economists think, for example, that monetary shocks drove the

U.S. Great Depression, but these economists disagree about the details of

the driving mechanism. Our analysis suggests that models in which financial

frictions show up primarily as investment wedges are not promising while

models in which financial frictions show up as efficiency or labor wedges

may well be. Thus, we conclude that researchers interested in developing

models in which monetary shocks lead to the Great Depression should focus on

detailed models in which financial frictions manifest themselves as

efficiency and labor wedges.



Other economists, including Cole and Ohanian (1999 and 2004) and Prescott

(1999), emphasize nonmonetary factors behind the Great Depression,

downplaying the importance of money and banking shocks. For such economists,

our analysis guides them to promising models, like that of Cole and Ohanian

(2004), in which fluctuations in the power of unions and cartels lead to

labor wedges, and other models in which poor government policies lead to

efficiency wedges.



In terms of method, the equivalence result provides the logical foundation

for the way our accounting procedure uses the measured wedges. At a

mechanical level, the wedges represent deviations in the prototype model's

first-order conditions and in its relationship between inputs and outputs.

One interpretation of these deviations, of course, is that they are simply

errors, so that their size indicates the goodness-of-fit of the model. Under

that interpretation, however, feeding the measured wedges back into the

model makes no sense. Our equivalence result leads to a more economically

useful interpretation of the deviations by linking them directly to classes

of models; that link provides the rationale for feeding the measured wedges

back into the model.



Also in terms of method, the accounting procedure goes beyond simply

plotting the wedges. Such plots, by themselves, are not useful in evaluating

the quantitative importance of competing mechanisms of business cycles

because they tell us little about the equilibrium responses to the wedges.

Feeding the measured wedges back into the prototype model and measuring the

model's resulting equilibrium responses is what allows us to discriminate

between competing mechanisms.



Finally, in terms of method, our decomposition of business cycle

fluctuations is quite different from traditional decompositions. Those

decompositions attempt to isolate the effects of (so-called) primitive

shocks on equilibrium outcomes by making identifying assumptions, typically

zero-one restrictions on variables and shocks. The problem with the

traditional approach is that finding identifying assumptions that apply to a

broad class of detailed models is hard. Hence, this approach is not useful

in pointing researchers toward classes of promising models. Our approach, in

contrast, can be applied to a broad class of detailed models. Our

equivalence results, which provide a mapping from wedges to frictions in

particular detailed models, play the role of the identifying assumptions in

the traditional approach. This mapping is detailed-model specific and is the

key to interpreting the properties of the wedges we document. For any

detailed model of interest, researchers can use the mapping that is relevant

for their model to learn whether it is promising. In this sense our

approach, while being purposefully less ambitious than the traditional

approach, is much more flexible than that approach.



Our accounting procedure is intended to be a useful first step in guiding

the construction of detailed models with various frictions, to help

researchers decide which frictions are quantitatively important to business

cycle fluctuations. The procedure is not a way to test particular detailed

models. If a detailed model is at hand, then it makes sense to confront that

model directly with the data. Nevertheless, our procedure is useful in

analyzing models with many frictions. For example, some researchers, such as

Bernanke, Gertler, and Gilchrist (1999) and Christiano, Gust, and Roldos

(2004), have argued that the data are well accounted for by models which

include a host of frictions (such as credit market frictions, sticky wages,

and sticky prices). Our analysis suggests that the features of these models

which primarily lead to investment wedges can be dropped while only modestly

affecting the models' ability to account for the data.



Our work here is related to a vast business cycle literature that we discuss

in detail after we describe and apply our new method.



\bigskip



\begin{center}

1. \textsc{Demonstrating the Equivalence Result}

\end{center}



Here we show how various detailed models with underlying distortions are

equivalent to a prototype growth model with one or more wedges.



\bigskip



\begin{center}

1.1. \emph{The Benchmark Prototype Economy}

\end{center}



The \emph{benchmark }prototype economy that we use later in our accounting

procedure is a stochastic growth model. In each period $t,$ the economy

experiences one of finitely many events $s_{t},$ which index the shocks. We

denote by $s^{t}=(s_{0},...,s_{t})$ the history of events up through and

including period $t$ and often refer to $s^{t}$ as the \textit{state}. The

probability, as of period $0$, of any particular history $s^{t}$ is $\pi

_{t}(s^{t})$. The initial realization $s_{0}$ is given. The economy has four

exogenous stochastic variables, all of which are functions of the underlying

random variable $s^{t}$: the \textit{efficiency wedge }$A_{t}(s^{t}),$ the 

\textit{labor wedge }$1-\tau _{lt}(s^{t}),$ the \textit{investment wedge} $%

1/[1+\tau _{xt}(s^{t})]$, and the \textit{government consumption wedge }$%

g_{t}(s^{t})$.



In the model, consumers maximize expected utility over per capita

consumption $c_{t}$ and per capita labor $l_{t}$, 

\[

\sum_{t=0}^{\infty }\sum_{s^{t}}\beta ^{t}\pi _{t}(s^{t})U\text{{\large (}}%

c_{t}(s^{t}),l_{t}(s^{t})\text{{\large )}}N_{t}, 

\]%

subject to the budget constraint 

\[

c_{t}+[1+\tau _{xt}(s^{t})]x_{t}(s^{t})=[1-\tau

_{lt}(s^{t})]w_{t}(s^{t})l_{t}(s^{t})+r_{t}(s^{t})k_{t}(s^{t-1})+T_{t}(s^{t}) 

\]%

and the capital accumulation law 

\begin{equation}

(1+\gamma _{n})k_{t+1}(s^{t})=(1-\delta )k_{t}(s^{t-1})+x_{t}(s^{t}),

\label{1}

\end{equation}%

where $k_{t}(s^{t-1})$ denotes the per capita capital stock, $x_{t}(s^{t})$

per capita investment, $w_{t}(s^{t})$ the wage rate, $r_{t}(s^{t})$ the

rental rate on capital, $\beta $ the discount factor, $\delta $ the

depreciation rate of capital, $N_{t}$ the population with growth rate equal

to $1+\gamma _{n}$, and $T_{t}(s^{t})$ per capita lump-sum transfers.



The production function is $F${\large (}$k_{t}(s^{t-1}),(1+\gamma

)^{t}l_{t}(s^{t})${\large )}, where $1+\gamma $ is the rate of

labor-augmenting technical progress, which is assumed to be a constant.

Firms maximize profits given by $A_{t}(s^{t})F${\large (}$%

k_{t}(s^{t-1}),(1+\gamma )^{t}l_{t}(s^{t})${\large )}$%

-r_{t}(s^{t})k_{t}(s^{t-1})-w_{t}(s^{t})l_{t}(s^{t})$.



The equilibrium of this benchmark prototype economy is summarized by the

resource constraint, 

\begin{equation}

c_{t}(s^{t})+x_{t}(s^{t})+g_{t}(s^{t})=y_{t}(s^{t}),  \label{resource}

\end{equation}%

where $y_{t}(s^{t})$\ denotes per capita output, together with 

\begin{equation}

y_{t}(s^{t})=A_{t}(s^{t})F\text{{\large (}}k_{t}(s^{t-1}),(1+\gamma

)^{t}l_{t}(s^{t})\text{{\large )}},  \label{AA}

\end{equation}%

\begin{equation}

-\frac{U_{lt}(s^{t})}{U_{ct}(s^{t})}=[1-\tau

_{lt}(s^{t})]A_{t}(s^{t})(1+\gamma )^{t}F_{lt},\text{ and }  \label{taun}

\end{equation}%

\begin{equation}

U_{ct}(s^{t})[1+\tau _{xt}(s^{t})]  \label{taux}

\end{equation}%

\[

=\beta \sum_{s^{t+1}}\pi

_{t}(s^{t+1}|s^{t})U_{ct+1}(s^{t+1})\{A_{t+1}(s^{t+1})F_{kt+1}(s^{t+1})+(1-%

\delta )[1+\tau _{xt+1}(s^{t+1})]\}, 

\]%

where, here and throughout, notations like $U_{ct}$, $U_{lt}$, $F_{lt}$, and 

$F_{kt}$ denote the derivatives of the utility function and the production

function with respect to their arguments and $\pi _{t}(s^{t+1}|s^{t})$

denotes the conditional probability $\pi _{t}(s^{t+1})/\pi _{t}(s^{t})$. We

assume that $g_{t}(s^{t})$ fluctuates around a trend of $(1+\gamma )^{t}.$



Notice that in this benchmark prototype economy, the efficiency wedge

resembles a blueprint technology parameter, and the labor wedge and the

investment wedge resemble tax rates on labor income and investment. Other

more elaborate models could be considered, models with other kinds of

frictions that look like taxes on consumption or on capital income.

Consumption taxes induce a wedge between the consumption-leisure marginal

rate of substitution and the marginal product of labor in the same way as do

labor income taxes. Such taxes, if time-varying, also distort the

intertemporal margins in (\ref{taux}). Capital income taxes induce a wedge

between the intertemporal marginal rate of substitution and the marginal

product of capital which is only slightly different from the distortion

induced by a tax on investment. We experimented with intertemporal

distortions that resemble capital income taxes rather than investment taxes

and found that our substantive conclusions are unaffected. (For details, see

Chari, Kehoe, and McGrattan (2006), hereafter referred to as the \emph{%

technical appendix}.)



We emphasize that each of the wedges represents the overall distortion to

the relevant equilibrium condition of the model. For example, distortions

both to labor supply affecting consumers and to labor demand affecting firms

distort the static first-order condition (\ref{taun}). Our labor wedge

represents the sum of these distortions. Thus, our method identifies the

overall wedge induced by both distortions and does not identify each

separately. Likewise, liquidity constraints on consumers distort the

consumer's intertemporal Euler equation, while investment financing

frictions on firms distort the firm's intertemporal Euler equation. Our

method combines the Euler equations for the consumer and the firm and

therefore identifies only the overall wedge in the combined Euler equation

given by (\ref{taux}). We focus on the overall wedges because what matters

in determining business cycle fluctuations is the overall wedges, not each

distortion separately.



\bigskip



\begin{center}

1.2. \emph{The Mapping---From Frictions to Wedges}

\end{center}



Now we illustrate the mapping between detailed economies and prototype

economies for two types of wedges. We show that input-financing frictions in

a detailed economy map into efficiency wedges in our prototype economy.

Sticky wages in a monetary economy map into our prototype (real) economy

with labor wedges. In an appendix, we show as well that investment-financing

frictions map into investment wedges and that fluctuations in net exports in

an open economy map into government consumption wedges in our prototype

(closed) economy. In general, our approach is to show that the frictions

associated with specific economic environments manifest themselves as

distortions in first-order conditions and resource constraints in a growth

model. We refer to these distortions as \emph{wedges}.



We choose simple models in order to illustrate how the detailed models map

into the prototypes. Since many models map into the same configuration of

wedges, identifying one particular configuration does not uniquely identify

a model; rather, it identifies a whole class of models consistent with that

configuration. In this sense, our method does not uniquely determine the

model most promising to analyze business cycle fluctuations. It does,

however, guide researchers to focus on the key margins that need to be

distorted in order to capture the nature of the fluctuations.



\bigskip \newpage



\noindent a. \emph{Efficiency Wedges}



In many economies, underlying frictions either within or across firms cause

factor inputs to be used inefficiently. These frictions in an underlying

economy often show up as aggregate productivity shocks in a prototype

economy similar to our benchmark economy. Schmitz (2005) presents an

interesting example of within-firm frictions resulting from work rules that

lower measured productivity at the firm level. Lagos (2006) studies how

labor market policies lead to misallocations of labor across firms and,

thus, to lower aggregate productivity. And Chu (2001) and Restuccia and

Rogerson (2003) show how government policies at the levels of plants and

establishments lead to lower aggregate productivity.



Here we develop a detailed economy with input-financing frictions and use it

to make two points. This economy illustrates the general idea that frictions

which lead to inefficient factor utilization map into efficiency wedges in a

prototype economy. Beyond that, however, the economy also demonstrates that

financial frictions can show up as efficiency wedges rather than as

investment wedges. In our detailed economy, financing frictions lead some

firms to pay higher interest rates for working capital than do other firms.

Thus, these frictions lead to an inefficient allocation of inputs across

firms.



\bigskip



\noindent $\square $ \emph{A Detailed Economy With Input-Financing Frictions}



Consider a simple detailed economy with financing frictions which distort

the allocation of intermediate inputs across two types of firms. Both types

of firms must borrow to pay for an intermediate input in advance of

production. One type of firm is more financially constrained, in the sense

that it pays a higher interest rate on borrowing than does the other type.

We think of these frictions as capturing the idea that some firms, such as

small firms, often have difficulty borrowing. One motivation for the higher

interest rate faced by the financially constrained firms is that moral

hazard problems are more severe for small firms.



Specifically, consider the following economy. Aggregate gross output $q_{t}$

is a combination of the gross output $q_{it}$ from the economy's two

sectors, indexed $i=1,2$, where 1 indicates the sector of firms that are

more financially constrained and 2 the sector of firms that are less

financially constrained. The sectors' gross output is combined\ according to 

\begin{equation}

q_{t}=q_{1t}^{\phi }q_{2t}^{1-\phi },  \label{q}

\end{equation}

where $0<\phi <1.$ The representative producer of the gross output $q_{t}$

chooses $q_{1t}$ and $q_{2t}$ to solve this problem: 

\[

\max \ q_{t}-p_{1t}q_{1t}-p_{2t}q_{2t} 

\]

subject to (\ref{q}), where $p_{it}$ is the price of the output of sector $i$%

.



The resource constraint for gross output in this economy is 

\begin{equation}

c_{t}+k_{t+1}+m_{1t}+m_{2t}=q_{t}+(1-\delta )k_{t},  \label{gross output}

\end{equation}

where $c_{t}$ is consumption, $k_{t}$ is the capital stock, and $m_{1t}$ and 

$m_{2t}$ are intermediate goods used in sectors 1 and 2, respectively. Final

output, given by $y_{t}=$ $q_{t}-$ $m_{1t}-$ $m_{2t}$, is gross output less

the intermediate goods used.



The gross output of each sector $i$, $q_{it}$, is made from intermediate

goods $m_{it}$ and a composite value-added good $z_{it}$ according to 

\begin{equation}

q_{it}=m_{it}^{\theta }z_{it}^{1-\theta },  \label{stupid}

\end{equation}

where $0<\theta <1.$ The composite value-added good is produced from capital 

$k_{t}$ and labor $l_{t}$ according to 

\begin{equation}

z_{1t}+z_{2t}=z_{t}=F(k_{t},l_{t}).  \label{z}

\end{equation}



The producer of gross output of sector $i$ chooses the composite good $%

z_{it} $ and the intermediate good $m_{it}$ to solve this problem: 

\[

\max \ p_{it}q_{it}-v_{t}z_{it}-R_{it}m_{it} 

\]

subject to (\ref{stupid}). Here $v_{t}$ is the price of the composite good

and $R_{it}$ is the gross within-period interest rate paid on borrowing by

firms in sector $i.$ If firms in sector $1$ are more financially constrained

than those in sector 2, then $R_{1t}>R_{2t}.$ Let $R_{it}=R_{t}(1+\tau

_{it}) $, where $R_{t}$ is the rate consumers earn within period $t$ and $%

\tau _{it} $ measures the within-period spread, induced by financing

constraints, between the rate paid to consumers who save and the rate paid

by firms in sector $i$. Since consumers do not discount utility within the

period, $R_{t}=1.$



In this economy, the representative producer of the composite good $z_{t}$

chooses $k_{t}$ and $l_{t}$ to solve this problem: 

\[

\max \ v_{t}z_{t}-w_{t}l_{t}-r_{t}k_{t} 

\]

subject to (\ref{z}), where $w_{t}$ is the wage rate and $r_{t}$ is the

rental rate on capital.



Consumers solve this problem: 

\begin{equation}

\max \ \sum_{t=0}^{\infty }\beta ^{t}U(c_{t},l_{t})  \label{consumer}

\end{equation}%

subject to 

\[

c_{t}+k_{t+1}=r_{t}k_{t}+w_{t}l_{t}+(1-\delta )k_{t}+T_{t}, 

\]%

where $l_{t}=l_{1t}+l_{2t}$ is the economy's total labor supply and $%

T_{t}=R_{t}\sum_{i}\tau _{it}m_{it}$ lump-sum transfers. Here we assume that

the financing frictions act like distorting taxes, and the proceeds are

rebated to consumers. If, instead, we assumed that these frictions

represent, say, lost gross output, then we would adjust the economy's

resource constraint (\ref{gross output}) appropriately.



\bigskip



\noindent $\square $ \emph{The Associated Prototype Economy With Efficiency

Wedges}



Now consider a version of the benchmark prototype economy that will have the

same aggregate allocations as the input-financing frictions economy just

detailed. This prototype economy is identical to our benchmark prototype

except that the new prototype economy has an investment wedge that resembles

a tax on capital income rather than a tax on investment. Here the government

consumption wedge is set equal to zero.



Now the consumer's budget constraint is 

\begin{equation}

c_{t}+k_{t+1}=(1-\tau _{kt})r_{t}k_{t}+(1-\tau _{lt})w_{t}l_{t}+(1-\delta

)k_{t}+T_{t},  \label{budget A}

\end{equation}

and the efficiency wedge is 

\begin{equation}

A_{t}=\kappa (a_{1t}^{1-\phi }a_{2t}^{\phi })^{\frac{\theta }{1-\theta }%

}[1-\theta (a_{1t}+a_{2t})],  \label{Ainput}

\end{equation}

where $a_{1t}=\phi /(1+\tau _{1t})$, $a_{2t}=(1-\phi )/(1+\tau _{2t}),$ $%

\kappa =[\phi ^{\phi }(1-\phi )^{1-\phi }\theta ^{\theta }]^{\frac{1}{%

1-\theta }},$ and $\tau _{1t}$ and $\tau _{2t}$ are the interest rate

spreads in the detailed economy.



Comparing the first-order conditions in the detailed economy with

input-financing frictions to those of the associated prototype economy with

efficiency wedges leads immediately to this proposition: \bigskip



\textsc{Proposition }1:\textit{\ }\emph{Consider the prototype economy with

resource constraint} (\ref{resource}) \emph{and consumer budget constraint} (%

\ref{budget A}) \emph{with exogenous processes for the efficiency wedge }$%

A_{t}$\emph{\ given in }(\ref{Ainput}), \emph{the labor wedge given by } 

\begin{equation}

\frac{1}{1-\tau _{lt}}=\frac{1}{1-\theta }\left[ 1-\theta \left( \frac{\phi 

}{1+\tau _{1t}^{\ast }}+\frac{1-\phi }{1+\tau _{2t}^{\ast }}\right) \right] ,

\label{tax}

\end{equation}

\emph{and the investment wedge given by} $\tau _{kt}=\tau _{lt}$\emph{,} 

\emph{where} $\tau _{1t}^{\ast }$ \emph{and} $\tau _{2t}^{\ast }$ \emph{are

the interest rate spreads from the detailed economy with input-financing

frictions. Then the equilibrium allocations for aggregate variables in the

detailed economy are equilibrium allocations in this prototype economy.}%

\bigskip



Consider the following special case of Proposition 1 in which only the

efficiency wedge fluctuates. Specifically, suppose that in the detailed

economy the interest rate spreads $\tau _{1t}$ and $\tau _{2t}$ fluctuate

over time, but in such a way that the weighted average of these spreads, 

\begin{equation}

a_{1t}+a_{2t}=\frac{\phi }{1+\tau _{1t}}+\frac{1-\phi }{1+\tau _{2t}},

\label{average wedge}

\end{equation}

is constant while $a_{1t}^{1-\phi }a_{2t}^{\phi }$ fluctuates. Then from (%

\ref{tax}) we see that the labor and investment wedges are constant, and

from (\ref{Ainput}) we see that the efficiency wedge fluctuates. In this

case, on average, financing frictions are unchanged, but relative

distortions fluctuate. An outside observer who attempted to fit the data

generated by the detailed economy with input-financing frictions to the

prototype economy would identify the fluctuations in relative distortions

with fluctuations in technology and would see no fluctuations in either the

labor wedge $1-\tau _{lt}$ or the investment wedge $\tau _{kt}$. In

particular, periods in which the relative distortions increase would be

misinterpreted as periods of technological regress.



\bigskip



\noindent b. \emph{Labor Wedges}



Now we show that a monetary economy with sticky wages is equivalent to a

(real) prototype economy with labor wedges. In the detailed economy, the

shocks are to monetary policy, while in the prototype economy, the shocks

are to the labor wedge.



\bigskip



\noindent $\square $ \emph{A Detailed Economy With Sticky Wages}



Consider a monetary economy populated by a large number of identical,

infinitely lived consumers. The economy consists of a competitive final

goods producer and a continuum of monopolistically competitive unions that

set their nominal wages in advance of the realization of shocks to the

economy. Each union represents all consumers who supply a specific type of

labor.



In each period $t,$ the commodities in this economy are a

consumption-capital good, money, and a continuum of differentiated types of

labor, indexed by $j$ $\in \lbrack 0,1]$. The technology for producing final

goods from capital and a labor aggregate at history, or state, $s^{t}$ has

constant returns to scale and is given by $y(s^{t})=F${\large (}$%

k(s^{t-1}),l(s^{t})${\large )}$,$ where $y(s^{t})$ is output of the final

good, $k(s^{t-1})$ is capital, and 

\begin{equation}

l(s^{t})=\left[ \int l(j,s^{t})^{v}\,dj\right] ^{\frac{1}{v}}  \label{labor}

\end{equation}

is an aggregate of the differentiated types of labor $l(j,s^{t})$.



The final goods producer in this economy behaves competitively. This

producer has some initial capital stock $k(s^{-1})$ and accumulates capital

according to $k(s^{t})=(1-\delta )k(s^{t-1})+x(s^{t})$, where $x(s^{t})$ is

investment. The present discounted value of profits for this producer is 

\begin{equation}

\ \sum_{t=0}^{\infty }\sum_{s^{t}}Q(s^{t})\left[

P(s^{t})y(s^{t})-P(s^{t})x(s^{t})-W(s^{t-1})l(s^{t})\right] ,

\label{profits}

\end{equation}

where $Q(s^{t})$ is the price of a dollar at $s^{t}$ in an abstract unit of

account, $P(s^{t})$ is the dollar price of final goods at $s^{t}$, and $%

W(s^{t-1})$ is the aggregate nominal wage at $s^{t}$ which depends on only $%

s^{t-1}$ because of wage stickiness.



The producer's problem can be stated in two parts. First, the producer

chooses sequences for capital $k(s^{t-1}),$ investment $x(s^{t}),$ and

aggregate labor $l(s^{t})$ in order to maximize (\ref{profits}) given the

production function and the capital accumulation law. The first-order

conditions can be summarized by 

\begin{equation}

P(s^{t})F_{l}(s^{t})=W(s^{t-1})\text{ and}  \label{foc 1}

\end{equation}

\begin{equation}

Q(s^{t})P(s^{t})=\sum_{s_{t+1}}Q(s^{t+1})P(s^{t+1})[F_{k}(s^{t+1})+1-\delta

].  \label{foc 2}

\end{equation}

Second, for any given amount of aggregate labor $l(s^{t}),$ the producer's

demand for each type of differentiated labor is given by the solution to



\begin{equation}

\min_{\{l(j,s^{t})\},j\in \lbrack 0,1]}\int \,W(j,s^{t-1})l(j,s^{t})\,dj

\label{wages}

\end{equation}

subject to (\ref{labor}); here $W(j,s^{t-1})$ is the nominal wage for

differentiated labor of type $j$. Nominal wages are set by unions before the

realization of the event in period $t$; thus, wages depend on, at most, $%

s^{t-1}.$ The demand for labor of type $j$ by the final goods producer is 

\begin{equation}

l^{d}(j,s^{t})=\left[ \frac{W(s^{t-1})}{W(j,s^{t-1})}\right] ^{\frac{1}{1-v}%

}l(s^{t}),  \label{labor demand}

\end{equation}

where $W(s^{t-1})\equiv \left[ \int \,W(j,s^{t-1})^{\frac{v}{v-1}}\,dj\right]

^{\frac{v-1}{v}}$ is the aggregate nominal wage. The minimized value in (\ref%

{wages}) is, thus, $W(s^{t-1})l(s^{t})$.



In this economy, consumers can be thought of as being organized into a

continuum of unions indexed by $j$. Each union consists of all the consumers

in the economy with labor of type $j$. Each union realizes that it faces a

downward-sloping demand for its type of labor, given by (\ref{labor demand}%

). In each period, the new wages are set before the realization of the

economy's current shocks.



The preferences of a representative consumer in the $j$th union is 

\begin{equation}

\sum_{t=0}^{\infty }\sum_{s^{t}}\,\beta ^{t}\pi _{t}(s^{t})\,\text{{\large [}%

}U\text{{\large (}}c(j,s^{t}),l(j,s^{t})\text{{\large )}}+V\text{{\large (}}%

M(j,s^{t})/P(s^{t})\text{{\large )]}},  \label{preferences}

\end{equation}%

where $c(j,s^{t}),l(j,s^{t}),M(j,s^{t})$ are the consumption, labor supply,

and money holdings of this consumer, and $P(s^{t})$ is the economy's overall

price level. Note that the utility function is separable in real balances.

This economy has complete markets for state-contingent nominal claims. The

asset structure is represented by a set of complete, contingent, one-period

nominal bonds. Let $B(j,s^{t+1})$ denote the consumers' holdings of such a

bond purchased in period $t$ at history $s^{t}$, with payoffs contingent on

some particular event $s_{t+1}$ in $t+1,$ where $s^{t+1}=(s^{t},s_{t+1})$.

One unit of this bond pays one dollar in period $t+1$ if the particular

event $s_{t+1}$ occurs and $0$ otherwise. Let $Q(s^{t+1}|s^{t})$ denote the

dollar price of this bond in period $t$ at history $s^{t}$, where $%

Q(s^{t+1}|s^{t})=Q(s^{t+1})/Q(s^{t}).$



The problem of the $j$th union is to maximize (\ref{preferences}) subject to

the budget constraint 

\begin{eqnarray*}

&&P(s^{t})c(j,s^{t})+M(j,s^{t})+\sum_{s_{t+1}}Q(s^{t+1}|s^{t})B(j,s^{t+1}) \\

&&\qquad \qquad \qquad \leq

W(j,s^{t-1})l(j,s^{t})+M(j,s^{t-1})+B(j,s^{t})+P(s^{t})T(s^{t})+D(s^{t}),

\end{eqnarray*}

the constraint $l(j,s^{t})=l^{d}(j,s^{t}),$ and the borrowing constraint $%

B(s^{t+1})\geq -P(s^{t})\overline{b}$, where $l^{d}(j,s^{t})$ is given by (%

\ref{labor demand})$.$ Here $T(s^{t})$ is transfers and the positive

constant $\overline{b}$ constrains the amount of real borrowing by the

union. Also, $D(s^{t})=$ $%

P(s^{t})y(s^{t})-P(s^{t})x(s^{t})-W(s^{t-1})l(s^{t})$ are the dividends paid

by the firms$.$ The initial conditions $M(j,s^{-1})$ and $B(j,s^{0})$ are

given and assumed to be the same for all $j.$ Notice that in this problem,

the union chooses the wage and agrees to supply whatever labor is demanded

at that wage.



The first-order conditions for this problem can be summarized by 

\begin{equation}

\frac{V_{m}(j,s^{t})}{P(s^{t})}-\frac{U_{c}(j,s^{t})}{P(s^{t})}+\beta

\sum_{s_{t+1}}\pi (s^{t+1}|s^{t})\frac{U_{c}(j,s^{t+1})}{P(s^{t+1})}=0,

\label{money foc}

\end{equation}%

\begin{equation}

Q(s^{t}|s^{t-1})=\beta \pi _{t}(s^{t}|s^{t-1})\frac{U_{c}(j,s^{t})}{%

U_{c}(j,s^{t-1})}\frac{P(s^{t-1})}{P(s^{t})}\text{, and }  \label{Q}

\end{equation}%

\begin{equation}

W(j,s^{t-1})=-\frac{%

\sum_{s^{t}}Q(s^{t})P(s^{t})U_{l}(j,s^{t})/U_{c}(j,s^{t})l^{d}(j,s^{t})}{%

v\sum_{s^{t}}Q(s^{t})l^{d}(j,s^{t})}.  \label{wage setting 1}

\end{equation}%

Here $\pi _{t}(s^{t+1}|s^{t})=$ $\pi _{t}(s^{t+1})/\pi _{t}(s^{t})$ is the

conditional probability of $s^{t+1}$ given $s^{t}$. Notice that in a steady

state, ($\ref{wage setting 1})$ reduces to $W/P=(1/v)(-U_{l}/U_{c}),$ so

that real wages are set as a markup over the marginal rate of substitution

between labor and consumption. Given the symmetry among the unions, all of

them choose the same consumption, labor, money balances, bond holdings, and

wages, which are denoted simply by $c(s^{t}),$ $l(s^{t}),$ $M(s^{t}),$ $%

B(s^{t+1})$, and $W(s^{t}).$



Consider next the specification of the money supply process and the

market-clearing conditions for this sticky-wage economy. The nominal money

supply process is given by $M(s^{t})=\mu (s^{t})M(s^{t-1})$, where $\mu

(s^{t})$ is a stochastic process. New money balances are distributed to

consumers in a lump-sum fashion by having nominal transfers satisfy $%

P(s^{t})T(s^{t})=M(s^{t})-M(s^{t-1})$. The resource constraint for this

economy is $c(s^{t})+k(s^{t})=y(s^{t})+(1-\delta )k(s^{t-1}).$ Bond

market--clearing requires that $B(s^{t+1})=0.$



\bigskip



\noindent $\square $ \emph{The Associated Prototype Economy With Labor Wedges%

}



Consider now a real prototype economy with labor wedges and the production

function for final goods given above in the detailed economy with sticky

wages$.$ The representative firm maximizes (\ref{profits}) subject to the

capital accumulation law given above. The first-order conditions can be

summarized by (\ref{foc 1}) and (\ref{foc 2}). The representative consumer

maximizes 

\[

\sum_{t=0}^{\infty }\sum_{s^{t}}\,\beta ^{t}\pi _{t}(s^{t})\,U\text{{\large (%

}}c(s^{t}),l(s^{t})\text{{\large )}} 

\]%

subject to the budget constraint 

\[

c(s^{t})+\sum_{s_{t+1}}q(s^{t+1}|s^{t})b(s^{t+1})\leq {\large [}1-\tau

_{l}(s^{t}){\large ]}w(s^{t})l(s^{t})+b(s^{t})+v(s^{t})+d(s^{t}) 

\]%

with $w(s^{t})$ replacing $W(s^{t-1})/P(s^{t})$ and $q(s^{t+1}/s^{t})$

replacing $Q(s^{t+1})P(^{st+1})/Q(s^{t})P(s^{t})$ and a bound on real bond

holdings, where the lowercase letters $q,b,w,v,$ and $d$ denote the real

values of bond prices, debt, wages, lump-sum transfers, and dividends. Here

the first-order condition for bonds is identical to that in (\ref{Q}) once

symmetry has been imposed with $q(s^{t}/s^{t-1})$ replacing $%

Q(s^{t}/s^{t-1})P(s^{t})/P(s^{t-1})$. The first-order condition for labor is

given by 

\[

-\frac{U_{l}(s^{t})}{U_{c}(s^{t})}={\large [}1-\tau _{l}(s^{t}){\large ]}%

w(s^{t}). 

\]



Consider an equilibrium of the sticky wage economy for some given stochastic

process $M^{\ast }(s^{t})$ on money supply. Denote all of the allocations

and prices in this equilibrium with asterisks. Then this proposition can be

easily established:



\bigskip



\textsc{Proposition} 2: \emph{Consider the prototype economy just described

with labor wedges given by } 

\begin{equation}

1-\tau _{l}(s^{t})=-\frac{U_{l}^{\ast }(s^{t})}{U_{c}^{\ast }(s^{t})}\frac{1%

}{F_{l}^{\ast }(s^{t})},  \label{labor wedge}

\end{equation}

\emph{where} $U_{l}^{\ast }(s^{t}),U_{c}^{\ast }(s^{t})$, \emph{and} $%

F_{l}^{\ast }(s^{t})$ \emph{are evaluated at the equilibrium of the sticky

wage economy and where real transfers are equal to the real value of

transfers in the sticky wage economy adjusted for the interest cost of

holding money. Then the equilibrium allocations and prices in the sticky

wage economy are the same as those in the prototype economy.}



\bigskip



The proof of this proposition is immediate from comparing the first-order

conditions, the budget constraints, and the resource constraints for the

prototype economy with labor wedges to those of the detailed economy with

sticky wages. The key idea is that distortions in the sticky-wage economy

between the marginal product of labor implicit in (\ref{wage setting 1}) and

the marginal rate of substitution between leisure and consumption are

perfectly captured by the labor wedges (\ref{labor wedge}) in the prototype

economy.



\bigskip \newpage



\begin{center}

2. \textsc{The Accounting Procedure}

\end{center}



Having established our equivalence result, we now describe our accounting

procedure at a conceptual level and discuss a Markovian implementation of it.



Our procedure is to conduct experiments that isolate the marginal effect of

each wedge as well as the marginal effects of combinations of these wedges

on aggregate variables. In the experiment in which we isolate the marginal

effect of the efficiency wedge, for example, we hold the other wedges fixed

at some constant values in all periods. In conducting this experiment, we

ensure that the probability distribution of the efficiency wedge coincides

with that in the prototype economy. In effect, we ensure that agents'

expectations of how the efficiency wedge will evolve are the same as in the

prototype economy. For each experiment, we compare the properties of the

resulting equilibria to those of the prototype economy. These comparisons,

together with our equivalence results, allow us to identify promising

classes of detailed economies.



\bigskip



\begin{center}

2.1 \emph{The Accounting Procedure at a Conceptual Level}

\end{center}



Suppose for now that the stochastic process $\pi _{t}(s^{t})$ and the

realizations of the state $s^{t}$ in some particular episode are known.

Recall that the prototype economy has one underlying (vector-valued) random

variable, the state $s^{t},$ which has a probability of $\pi _{t}(s^{t}).$

All of the other stochastic variables, including the four wedges---the

efficiency wedge $A_{t}(s^{t}),$ the labor wedge $1-\tau _{lt}(s^{t}),$ the

investment wedge $1/[1+\tau _{xt}(s^{t})]$, and the government consumption

wedge $g_{t}(s^{t})$---are simply functions of this random variable. Hence,

when the state $s^{t}$ is known, so are the wedges.



To evaluate the effects of just the efficiency wedge, for example, we

consider an economy, referred to as an \textit{efficiency wedge alone

economy,} with the same underlying state $s^{t}$ and probability $\pi

_{t}(s^{t})$ and the same function $A_{t}(s^{t})$ for the efficiency wedge

as in the prototype economy, but in which the other three wedges are set to

constants, in that $\tau _{lt}(s^{t})=\bar{\tau}_{l},\tau _{xt}(s^{t})=\bar{%

\tau}_{x},$ and $g_{t}(s^{t})=\bar{g}.$ Note that this construction ensures

that the probability distribution of the efficiency wedge in this economy is

identical to that in the prototype economy.



For the efficiency wedge alone economy, we then compute the equilibrium

outcomes associated with the realizations of the state $s^{t}$ in a

particular episode and compare these outcomes to those of the economy with

all four wedges. We find this comparison to be of particular interest

because in our applications, the realizations $s^{t}$ are such that the

economy with all four wedges exactly reproduces the data on output, labor,

investment, and consumption.



In a similar manner, we define the \textit{labor wedge alone economy}, the 

\textit{investment wedge alone economy}, and the \textit{government

consumption wedge alone economy, }as well as economies with a combination of

wedges such as the \textit{efficiency and labor wedge economy}.



\bigskip



\begin{center}

2.2 \emph{A Markovian Implementation}

\end{center}



So far we have described our procedure assuming that we know the stochastic

process $\pi _{t}(s^{t})$ and that we can observe the state $s^{t}.$ In

practice, of course, we need to either specify the stochastic process a

priori or use data to estimate it, and we need to uncover the state $s^{t}$

from the data. Here we describe a set of assumptions that makes these

efforts easy. Then we describe in detail the three steps involved in

implementing our procedure.



We assume that the state $s^{t}$ follows a Markov process of the form $\pi

(s_{t}|s_{t-1})$ and that the wedges in period $t$ can be used to uniquely

uncover the event $s_{t}$, in the sense that the mapping from the event $%

s_{t}$ to the wedges $(A_{t},\tau _{lt},\tau _{xt},g_{t})$ is one-to-one and

onto. Given this assumption, without loss of generality, let the underlying

event $s_{t}=(s_{At},s_{lt},s_{xt},s_{gt}),$ and let $A_{t}(s^{t})=s_{At},%

\tau _{lt}(s^{t})=s_{lt},\tau _{xt}(s^{t})=s_{xt},$ and $%

g_{t}(s^{t})=s_{gt}. $ Note that we have effectively assumed that agents use

only past wedges to forecast future wedges and that the wedges in period $t$

are sufficient statistics for the event in period $t.$



The first step in our procedure is to use data on $y_{t},l_{t},x_{t},$ and $%

g_{t}$ from an actual economy to estimate the parameters of the Markov

process $\pi (s_{t}|s_{t-1}).\ $We can do so using a variety of methods,

including the maximum likelihood procedure described below.



The second step in our procedure is to uncover the event $s_{t}$ by

measuring the realized wedges. We measure the government consumption wedge

directly from the data as the sum of government spending and net exports. To

obtain the values of the other three wedges, we use the data and the model's

decision rules. With $y_{t}^{d},l_{t}^{d},x_{t}^{d},g_{t}^{d}$, and $%

k_{0}^{d}$ denoting the data and $y(s_{t},k_{t}),\ l(s_{t},k_{t}),$ and$\

x(s_{t},k_{t})$ denoting the decision rules of the model, the realized wedge

series $s_{t}^{d}$ solves 

\begin{equation}

y_{t}^{d}=y(s_{t}^{d},k_{t}),~l_{t}^{d}=l(s_{t}^{d},k_{t}),\text{ and }%

x_{t}^{d}=x(s_{t}^{d},k_{t}),  \label{decision}

\end{equation}%

with $k_{t+1}=(1-\delta )k_{t}+x_{t}^{d},$ $k_{0}=k_{0}^{d}$, and $%

g_{t}=g_{t}^{d}.$ Note that we construct a series for the capital stock

using the capital accumulation law (\ref{1}), data on investment $x_{t}$,

and an initial choice of capital stock $k_{0}$. In effect, we solve for the

three unknown elements of the vector $s_{t}$ using the three equations (\ref%

{AA})--(\ref{taux}) and thereby uncover the state. We use the associated

values for the wedges in our experiments.



Note that the four wedges account for all of the movement in output, labor,

investment, and government consumption, in that if we feed the four wedges

into the three decision rules in (\ref{decision}) and use $%

g_{t}(s_{t}^{d})=s_{gt}\ $along with the law of motion for capital, we

simply recover the original data.



Note also that, in measuring the realized wedges, the estimated stochastic

process plays a role in measuring only the investment wedge. To see that the

stochastic process does not play a role in measuring the efficiency and

labor wedges, note that these wedges can equivalently be directly calculated

from (\ref{AA}) and (\ref{taun}) without computing the equilibrium of the

model. In contrast, calculating the investment wedge requires computing the

equilibrium of the model because the right side of (\ref{taux}) has

expectations over future values of consumption, the capital stock, the

wedges, and so on. The equilibrium of the model depends on these

expectations and, therefore, on the stochastic process driving the wedges.



The third step in our procedure is to conduct experiments to isolate the

marginal effects of the wedges. To do that, we allow a subset of the wedges

to fluctuate as they do in the data while the others are set to constants.

To evaluate the effects of the efficiency wedge, we compute the decision

rules for the efficiency wedge alone economy, denoted $%

y^{e}(s_{t},k_{t}),l^{e}(s_{t},k_{t}),$ and $x^{e}(s_{t},k_{t}),$ in which  $%

A_{t}(s^{t})=s_{At},\tau _{lt}(s^{t})=\bar{\tau}_{l},\tau _{xt}(s^{t})=\bar{%

\tau}_{x},$ and $g_{t}(s^{t})=\bar{g}.$ Starting from $k_{0}^{d}$, we then

use $s_{t}^{d}$, the decision rules, and the capital accumulation law to

compute the realized sequence of output, labor, and investment, $%

y_{t}^{e},l_{t}^{e},$ and $x_{t}^{e}$, which we call the \textit{efficiency

wedge components} of output, labor, and investment. We compare these

components to output, labor, and investment in the data. Other components

are computed and compared similarly.



Notice that in this experiment we computed the decision rules for an economy

in which only one wedge fluctuates and the others are set to be constants in

all events. The fluctuations in the one wedge are driven by fluctuations in

a 4 dimensional state $s_{t}.$



Notice also that our experiments are designed to separate out the direct

effect and the forecasting effect of fluctuations in wedges. As a wedge

fluctuates, it directly affects either budget constraints or resource

constraints. This fluctuation also affects the forecasts of that wedge as

well as of other wedges in the future. Our experiments are designed so that

when we hold a particular wedge constant, we eliminate the direct effect of

that wedge, but we retain its forecasting effect on the other wedges. By

doing so, we ensure that expectations of the fluctuating wedges are

identical to those in the prototype economy.



Here we focus on one simple way to specify the expectations of agents:

assume they simply use past values of the wedges to forecast future values.

An extension of our Markovian procedure is to use past endogenous variables,

such as output, investment, consumption, and perhaps even asset prices such

as stock market values, in addition to past wedges to forecast future

wedges. Another approach is to simply specify these expectations directly,

as we did in our earlier work (Chari, Kehoe, and McGrattan (2002)) and then

conduct a variety of experiments to determine how the results change as the

specification is changed.



\begin{center}

3. \textsc{Applying the Accounting Application}

\end{center}



Now we demonstrate how to apply our accounting procedure to two U.S.

business cycle episodes: the Great Depression and the postwar recession of

1982. We then extend our analysis to the entire postwar period. (In the

technical appendix, we describe in detail our data sources, parameter

choices, computational methods, and estimation procedures.)



\bigskip



\begin{center}

3.1. \emph{Details of the Application }

\end{center}



To apply our accounting procedure, we use functional forms and parameter

values familiar from the business cycle literature. We assume that the

production function has the form $F(k,l)=k^{\alpha }l^{1-\alpha }$ and the

utility function the form $U(c,l)=\log c+\psi \log (1-l).$ We choose the

capital share $\alpha =.35$ and the time allocation parameter $\psi =2.24.$

We choose the depreciation rate $\delta $, the discount factor $\beta $, and

growth rates $\gamma $ and $\gamma _{n}$\ so that, on an annualized basis,

depreciation is 4.64\%, the rate of time preference 3\%, the population

growth rate 1.5\%, and the growth of technology 1.6\%.



To estimate the stochastic process for the state, we first specify a vector

autoregressive AR(1) process for the event $%

s_{t}=(s_{At},s_{lt},s_{xt},s_{gt})$ of the form 

\begin{equation}

s_{t+1}=P_{0}+Ps_{t}+\varepsilon _{t+1},  \label{s process}

\end{equation}%

where the shock $\varepsilon _{t}$ is i.i.d. over time and is distributed

normally with mean zero and covariance matrix $V.$ To ensure that our

estimate of $V$ is positive semidefinite, we estimate the lower triangular

matrix $Q$, where $V=QQ^{\prime }.$ The matrix $Q$ has no structural

interpretation. (In section 5, we elaborate on the contrast between our

decomposition and more traditional decompositions which impose structural

interpretations on $Q$.)



We then use a standard maximum likelihood procedure to estimate the

parameters $P_{0},P$, and $V$ of the vector AR(1) process for the wedges. In

doing so, we use the log-linear decision rules of the prototype economy and

data on output, labor, investment, and the sum of government consumption and

net exports.



For our Great Depression experiments, we proceed as follows. We discretize

the process ($\ref{s process})$ and simulate the economy using nonlinear

decision rules from a finite-element method. We use nonlinear decision rules

in these experiments because the shocks are so large that, for a given

stochastic process, the linear decision rules are a poor approximation to

the nonlinear decision rules. Of course, we would rather have used the

nonlinear decision rules in estimating the parameters of the vector AR(1)

process. We do not do so because this exercise is computationally demanding.

Instead we experiment by varying the parameters of the vector AR(1) process

and find that our results are very similar across these experiments.



For our postwar experiments, we use the log-linear decision rules and the

continuous state process ($\ref{s process})$.



In order to implement our accounting procedure, we must first adjust the

data to make them consistent with the theory. In particular, we adjust the

U.S. data on output and its components to remove sales taxes and to add the

service flow for consumer durables. For the pre--World War II period, we

remove military compensation as well. We estimate separate sets of

parameters for the stochastic process for wedges (\ref{s process}) for each

of our two historical episodes. The other parameters are the same in the two

episodes. (See our technical appendix for our rationale for this decision.)

The stochastic process parameters for the Great Depression analysis are

estimated using annual data for 1901--40; those for analysis after World War

II, using quarterly data for 1959:1--2004:3. In the Great Depression

analysis, we impose the additional restriction that the covariance between

the shocks to the government consumption wedge and those to the other wedges

is zero. This restriction avoids having the large movements in government

consumption associated with World War I dominate the estimation of the

stochastic process.



Table I displays the resulting estimated values for the parameters of the

coefficient matrices, $P$ and $Q$, and the associated confidence bands for

our two historical data periods. The stochastic process (\ref{s process})

with these values will be used by agents in our economy to form their

expectations about future wedges.



\bigskip



\begin{center}

3.2. \emph{Findings}

\end{center}



Now we describe the results of applying our procedure to two historical U.S.

business cycle episodes. In the Great Depression, the efficiency and labor

wedges play a central role for all variables considered. In the 1982

recession, the efficiency wedge plays a central role for output and

investment while the labor wedge plays a central role for labor. The

government consumption wedge plays no role in either period. Most

strikingly, neither does the investment wedge.



In reporting our findings, we remove a trend of 1.6\% from output,

investment, and the government consumption wedge. Both output and labor are

normalized to equal 100 in the base periods: 1929 for the Great Depression

and 1979:1 for the 1982 recession. In both of these historical episodes,

investment (detrended) is divided by the base period level of output. Since

the government consumption component accounts for virtually none of the

fluctuations in output, labor, and investment, we discuss the government

consumption wedge and its components only in our technical appendix. Here we

focus primarily on the fluctuations due to the efficiency, labor, and

investment wedges.



\bigskip



\noindent a.\emph{\ }\textit{The Great Depression}



Our findings for the period 1929--39, which includes the Great Depression,

are displayed in Figures 1--4. In sum, we find that the efficiency and labor

wedges account for essentially all of the movements of output, labor, and

investment in the Depression period and that the investment wedge actually

drives output the wrong way.



In Figure 1, we display actual U.S. output along with the three measured

wedges for that period: the efficiency wedge $A,$ the labor wedge $(1-\tau

_{l}),$ and the investment wedge $1/(1+\tau _{x})$. We see that the

underlying distortions revealed by the three wedges have different patterns.

The distortions that manifest themselves as efficiency and labor wedges

become substantially worse between 1929 and 1933. By 1939, the efficiency

wedge has returned to the 1929 trend level, but the labor wedge has not.

Over the period, the investment wedge fluctuates, but investment decisions

are generally less distorted, in the sense that $\tau _{x}$ is smaller

between 1932 and 1939 than it is in 1929. Note that this investment wedge

pattern does not square with models of business cycles in which financial

frictions increase in downturns and decrease in recoveries.



In Figure 2, we plot the 1929--39 data for U.S. output, labor, and

investment along with the model's predictions for those variables when the

model includes just one wedge. In terms of the data, note that labor

declines $27\%$ from 1929 to 1933 and stays relatively low for the rest of

the decade. Investment also declines sharply from 1929 to 1933 but partially

recovers by the end of the decade. Interestingly, in an algebraic sense,

about half of output's $36\%$ fall from 1929 to 1933 is due to the decline

in investment.



In terms of the model, we start by assessing the separate contributions of

the three wedges.



Consider first the contribution of the efficiency wedge. In Figure 2, we see

that with this wedge alone, the model predicts that output declines less

than it actually does in the data and that it recovers more rapidly. For

example, by 1933, predicted output falls about $30\%$ while U.S. output

falls about $36\%$. Thus, the efficiency wedge accounts for over $80\%$ of

the decline of output in the data. By 1939, predicted output is only about $%

6\%$ below trend rather than the observed $22\%.$ As can also be seen in

Figure 2, the reason for this predicted rapid recovery is that the

efficiency wedge accounts for only a small part of the observed movements in

labor in the data. By 1933 the fall in predicted investment is similar but

somewhat greater than that in the data. It recovers faster, however.



Consider next the contributions of the labor wedge. In Figure 2, we see that

with this wedge alone, the model predicts output due to the labor wedge to

fall by 1933 a little less than half as much as output falls in the data: $%

16\%$ vs. $36\%.$ By 1939, however, the labor wedge model's predicted output

completely captures the slow recovery: it predicts output falling $21\%$,

approximately as much as output does that year in the data. This model

captures the slow output recovery because predicted labor due to the labor

wedge also captures the sluggishness in labor after 1933 remarkably well.

The associated prediction for investment is a decline, but not the actual

sharp decline from 1929 to 1933.



Summarizing Figure 2, we can say that the efficiency wedge accounts for over

three-quarters of output's downturn during the Great Depression but misses

its slow recovery, while the labor wedge accounts for about one-half of this

downturn and essentially all of the slow recovery.



Now consider the investment wedge. In Figure 3, we again plot the data for

output, labor, and investment, but this time along with the contributions to

those variables that the model predicts are due to the investment wedge

alone. This figure demonstrates that the investment wedge's contributions

completely miss the observed movements in all three variables. The

investment wedge actually leads output to rise by about $9\%$ by 1933.



Together, then, Figures 2 and 3 suggest that the efficiency and labor wedges

account for essentially all of the movements of output, labor, and

investment in the Depression period and that the investment wedge accounts

for almost none. This suggestion is confirmed by Figure 4. There we plot the

combined contribution from the efficiency, labor, and (insignificant)

government consumption wedges (labeled \emph{Model With No Investment Wedge}%

). As can be seen from the figure, essentially all of the fluctuations in

output, labor, and investment can be accounted for by movements in the

efficiency and labor wedges. For comparison, we also plot the combined

contribution due to the labor, investment, and government consumption wedges

(labeled \emph{Model With No Efficiency Wedge}). This combination does not

do well. In fact, comparing Figures 2 and 4, we see that the model with this

combination is further from the data than the model with the labor wedge

component alone.



One issue of possible concern with our findings about the role of the

investment wedge is that measuring it is subtler than measuring the other

wedges. Recall that the measurement of this wedge depends on the details of

the stochastic process governing the wedges, whereas the size of the other

wedges can be inferred from static equilibrium conditions. To address this

concern, we conduct an additional experiment intended to give the model with

no efficiency wedge the best chance of accounting for the data.



In this experiment, we choose the investment wedge to be as large as it

needs to be for investment in the model to be as close as possible to

investment in the data, and we set the other wedges to be constants.

Predictions of this model, which we call the \emph{Model With Maximum

Investment Wedge}, turn out to poorly match the behavior of consumption in

the data. For example, from 1929 to 1933, consumption in the model rises

more than $8\%$ relative to trend while consumption in the data declines

about $28\%.$ (For details, see the technical appendix.) We label this poor

performance the \textit{consumption anomaly} of the investment wedge model.



Altogether, these findings lead us to conclude that distortions which

manifest themselves primarily as investment wedges played essentially no

useful role in the U.S. Great Depression.\bigskip



\noindent b.\emph{\ The 1982 Recession}



Now we apply our accounting procedure to a more typical U.S. business cycle:

the recession of 1982. Here we get basically the same results as with the

earlier period: the efficiency and labor wedges play primary roles in the

business cycle fluctuations, and the investment wedge plays essentially none.



We start here as we did in the Great Depression analysis, by displaying

actual U.S. output over the entire business cycle period---here,

1979--85---along with the three measured wedges for that period. In Figure

5, we see that output falls nearly $10\%$ relative to trend between 1979 and

1982 and by 1985 is back up to about $1\%$ below trend. We also see that the

efficiency wedge falls between 1979 and 1982 and by 1985 is still a little

more than $3\%$ below trend. The labor wedge also worsens from 1979 to 1982,

but it improves substantially by 1985. The investment wedge, meanwhile,

fluctuates until 1983 and improves thereafter.



An analysis of the effects of the wedges separately for the 1979--85 period

is in Figures 6 and 7. In Figure 6, we see that the model with the

efficiency wedge alone produces a decline in output from 1979 to 1982 of $%

6\% $, which is about $60\%$ of the actual decline in that period. Here

output recovers a bit more slowly than in the data, but seems to otherwise

generally parallel the data's movements. The model with the labor wedge

alone produces a decline in output from 1979 to 1982 of only about $3\%.$ In

Figure 7, we see that the model with just the investment wedge produces

essentially no fluctuations in output.



Now we examine how well a combination of wedges reproduces the data for the

1982 recession period just as we did for the Depression period. In Figure 8,

we plot the movements in output, labor, and investment during 1979--85 due

to two combinations of wedges. One is the combined effects of the

efficiency, labor, and (insignificant) government consumption components

(labeled \emph{Model With No Investment Wedge}). In terms of output, this

combination mimics the decline in output until 1982 extremely well and

produces a slightly shallower recovery than in the data. The other is the

combination of the labor, investment, and government components (labeled 

\emph{Model With No Efficiency Wedge}), which produces a modest decline in

output relative to the data. In Figures 6, 7, and 8, we see clearly that in

the model with no efficiency wedge the labor wedge accounts for essentially

all of the decline and the investment wedge, essentially none. \bigskip



\begin{center}

3.3. \emph{Extending the Analysis to the Entire Postwar Period}

\end{center}



So far we have analyzed the wedges and their contributions for specific

episodes. The findings for both episodes suggest that frictions in detailed

models which manifest themselves as investment wedges in the benchmark

prototype economy play, at best, a tertiary role in accounting for business

cycle fluctuations. Do our findings apply beyond those particular episodes?

We attempt to extend our analysis to the entire postwar period by developing

some summary statistics for the period from 1959:1 through 2004:3 using

HP-filtered data. We first consider the standard deviations of the wedges

relative to output as well as correlations of the wedges with each other and

with output at various leads and lags. We then consider the standard

deviations and the cross correlations of output due to each wedge. These

statistics summarize salient features of the wedges and their role in output

fluctuations for the entire postwar sample. We think of the wedge statistics

as analogs of our plots of the wedges and the output statistics as analogs

of our plots of output due to just one wedge.\footnote{%

In Chari, Kehoe, and McGrattan (2004), we apply a spectral method to

determine the contributions of the wedges based on the population properties

of the stochastic process generated by the model. We do this for both

periods and find that the investment wedge plays only a modest role in both.}

The results suggest that our earlier findings do hold up, at least in a

relative sense: the investment wedge seems to play a larger role over the

entire postwar period than in the 1982 recession, but its effects are still

quite modest compared to those of the other wedges.



In Tables II and III, we display standard deviations and cross correlations

calculated using HP-filtered data for the postwar period. Panel A of Table

II shows that the efficiency, labor, and investment wedges are positively

correlated with output both contemporaneously and for several leads and

lags. In contrast, the government consumption wedge is somewhat negatively

correlated with output, both contemporaneously and for several leads and

lags. (Note that the government consumption wedge is the sum of government

consumption and net exports and that net exports are negatively correlated

with output.) Panel B of Table II shows that the cross correlations of the

efficiency, labor, and investment wedges are generally positive.



Table III summarizes various statistics of the movements of output over this

period due to each wedge. Consider panel A, and focus first on the output

fluctuations due to the efficiency wedge. Table III shows that output

movements due to this wedge have a standard deviation which is $73\%$ of

that of output in the data. These movements are highly positively correlated

with output in the data, both contemporaneously and for several leads and

lags. These statistics are consistent with our episodic analysis of the $%

1982 $ recession, which showed that the efficiency wedge can account for

about 60\% of the actual decline in output during that period and comoves

highly with it.



Consider next the role of the other wedges in the entire postwar period.

Return to Table III. In panel A, again, we see that output due to the labor

wedge alone fluctuates almost $60\%$ as much as does output in the data and

is positively correlated with it. Output due to the investment wedge alone

fluctuates less than a third as much as output in the data and is somewhat

positively correlated with it. Finally, output due to the government

consumption wedge alone fluctuates about $40\%$ as much as output in the

data and is somewhat negatively correlated with it. In panel B of Table III,

we see that output movements due to the efficiency and labor wedges as well

as the efficiency and investment wedges are positively correlated and that

the cross correlations of output movements due to the other wedges are

mostly essentially zero or negative.



All of our analyses using business cycle accounting thus seem to lead to the

same conclusion: to study business cycles, the most promising detailed

models to explore are those in which frictions manifest themselves primarily

as efficiency or labor wedges, not as investment wedges.



\bigskip



\begin{center}

4. \textsc{Interpreting Wedges With }\\[0pt]

\textsc{Alternative Technology or Preference Specifications}

\end{center}



In detailed economies with technology and preferences similar to those in

our benchmark prototype economy, the equivalence propositions proved thus

far provide a mapping between frictions in those detailed economies and

wedges in the prototype economy. Here we construct a similar mapping when

technology or preferences differ in the two types of economies. We then ask

if this alternative mapping changes our substantive conclusion that

financial frictions which manifest themselves primarily as investment wedges

are unlikely to play a primary role in accounting for business cycles. We

find that it does not.



When detailed economies have technology or preferences different from the

benchmark economy's, wedges in the benchmark economy can be viewed as

arising from two sources: frictions in the detailed economy and differences

in the specification of technology or preferences. While researchers could

simply use results from our benchmark prototype economy to draw inferences

about promising classes of models, drawing such inferences is easier with an

alternative approach. Basically, we decompose the wedges into their two

sources. To do that, construct an alternative prototype economy with

technology and preferences that do coincide with those in the detailed

economy, and repeat the business cycle accounting procedure with those two

economies. The part of the wedges in the benchmark prototype economy due to

frictions, then, will be the wedges in the alternative prototype economy,

while the remainder will be due to specification differences.



Here we use this approach to explore alternative prototype economies with

technology and preference specifications chosen because of their popularity

in the literature. These alternative specifications include variable instead

of fixed capital utilization, different labor supply elasticities, and

varying levels of costs to adjusting investment.



Two of these changes offer no help to investment wedges. Adding variable

capital utilization to the analysis shifts the relative contributions of the

efficiency and labor wedges to output's fluctuations---decreasing the

efficiency wedge's contribution and increasing the labor wedge's---but this

alternative specification leaves the investment wedge's contribution

definitely in third place. Adding different labor elasticities to the

analysis offers no help either.



The third specification change seems to give investment wedges a slightly

larger role, but still not a primary one. With investment adjustment costs

added to the analysis, the investment wedge in the benchmark prototype

economy depends on both the investment wedge and the marginal cost of

investment in the alternative prototype economy. We find that even if the

investment wedge is constant in the benchmark economy, it will worsen during

recessions and improve during booms in the alternative economy. With our

measured wedges, this finding suggests that with large enough adjustment

costs, investment wedges in detailed economies could play a significant role

in business cycle fluctuations.



To study this possibility, we investigate the effects of two parameter

values for adjustment costs: one at the level used by Bernanke, Gertler, and

Gilchrist (1999), the \emph{BGG level}, and one we consider \emph{extreme},

at four times that level. For the Great Depression period, we find that for

both adjustment cost levels, investment wedges play only a minor role. For

the 1982 recession period, we find that these wedges play a very small role

with BGG level costs and a somewhat larger but still modest role with the

much higher costs. These findings suggest that researchers who think

adjustment costs are extremely high may want to include in their models

financial frictions that manifest themselves as investment wedges. Such

models are not likely to do well, however, unless they also include other

frictions that play the primary role in business cycle fluctuations.



\bigskip



\begin{center}

4.1. \emph{Details of Alternative Specifications}

\end{center}



\noindent a. \emph{Variable Capital Utilization}



We begin with an extreme view about the amount of variability in capital

utilization.



Our specification of the technology which allows for variable capital

utilization follows the work of Kydland and Prescott (1988) and Hornstein

and Prescott (1993). We assume that the production function is now 

\begin{equation}

y=A(kh)^{\alpha }(nh)^{1-\alpha },  \label{variable y}

\end{equation}%

where $n$ is the number of workers employed and $h$ is the length (or \emph{%

hours}) of the workweek. The labor input is, then, $l=nh.$



In the data, we measure only the labor input $l$ and the capital stock $k$.

We do not directly measure $h$ or $n.$ The benchmark specification for the

production function can be interpreted as assuming that all of the observed

variation in measured labor input $l$ is in the number of workers and that

the workweek $h$ is constant$.$ Under this interpretation, our benchmark

specification with \emph{fixed capital utilization} correctly measures the

efficiency wedge (up to the constant $h).$



Now we investigate the opposite extreme: assume that the number of workers $%

n $ is constant and that all the variation in labor is from the workweek $h.$

Under this\textit{\ variable capital utilization} specification, the

services of capital $kh$ are proportional to the product of the stock $k$

and the labor input $l,$ so that variations in the labor input induce

variations in the flow of capital services. Thus, the capital utilization

rate is proportional to the labor input $l,$ and the efficiency wedge is

proportional to $y/k^{\alpha }.$



Consider an alternative prototype economy, denoted \emph{economy 2},

identical to a deterministic version of our benchmark prototype economy,

denoted \emph{economy 1}, except that the production function is now given

by $y=Ak^{\alpha }l.$ Let the sequence of wedges and the equilibrium

outcomes in the two economies be $(A_{it},\tau _{lit},\tau _{xit})$ and $%

(y_{it},c_{it},l_{it},x_{it})$ for $i=1,2.$ We then have this proposition:



\bigskip



\textsc{Proposition }3:\textit{\ }\emph{If the sequence of wedges for an

alternative prototype economy }2\emph{\ are related to the wedges in the

prototype economy }1\emph{\ by }$A_{2t}=A_{1t}l_{1t}^{-\alpha }$, $1-\tau

_{l2t}=$ $(1-\alpha )(1-\tau _{l1t}),$ \emph{and} $\tau _{x2t}=\tau _{x1t},$ 

\emph{then the equilibrium outcomes for the two economies coincide. }

\bigskip



\textsc{Proof:} We prove this proposition by showing that the equilibrium

conditions of economy $2$ are satisfied at the equilibrium outcomes of

economy $1.$ Since $y_{1t}=A_{1t}k_{1t}^{\alpha }l_{1t}^{1-\alpha },$ using

the definition of $A_{2t},$ we have that $y_{1t}=A_{2t}k_{1t}^{\alpha

}l_{1t}.$ The first-order condition for labor in economy $1$ is 

\[

-\frac{U_{lt}(c_{1t},l_{1t})}{U_{ct}(c_{1t},l_{1t})}=(1-\tau _{l1t})\frac{%

(1-\alpha )y_{1t}}{l_{1t}}. 

\]%

Using the definition of $\tau _{l2t},$ we have that 

\[

-\frac{U_{lt}(c_{1t},l_{1t})}{U_{ct}(c_{1t},l_{1t})}=(1-\tau _{l2t})\frac{%

y_{1t}}{l_{1t}}. 

\]%

The rest of the equations governing the equilibrium are unaffected. 

%TCIMACRO{\TeXButton{TeX field}{\hspace{1.75in}} }%

%BeginExpansion

\hspace{1.75in}

%EndExpansion

$\ \ Q.E.D.$



\bigskip



Note that even if the efficiency wedge in the alternative prototype economy

does not fluctuate, the associated efficiency wedge in the prototype economy

will. Proposition 3 also implies that if $\tau _{x1t}$ is a constant, so

that the contribution of the investment wedge to fluctuations in economy $1$

is zero, then $\tau _{x2t}$ is also a constant; hence, the contribution of

the investment wedge to fluctuations in economy $2$ is also zero. Extending

this proposition to a stochastic environment is immediate.



Now suppose that we are interested in detailed economies with variable

capital utilization. We use the alternative prototype economy to ask whether

this change affects our substantive conclusions. In answering this question,

we reestimate the parameters of the stochastic process for the underlying

state. (For details, see the technical appendix.)\ 



Variable capital utilization can induce significant changes in the measured

efficiency wedge. To see these changes, in Figure 9, we plot the measured

efficiency wedges for these two specifications of capital utilization during

the Great Depression period (with it \emph{fixed} in the benchmark economy

and \emph{variable} now). Clearly, when capital utilization is variable

rather than fixed, the efficiency wedge falls less and recovers more by

1939. (For the other wedges, see the technical appendix).



In Figure 10, we plot the data and the predicted output due to the

efficiency and labor wedges for the 1930s when the model includes variable

capital utilization. Comparing Figures 10 and 2, we see that with the

remeasured efficiency wedge, the labor wedge plays a much larger role in

accounting for the output downturn and slow recovery and\ the efficiency

wedge plays a much smaller role.



In Figure 11, we plot the three data series again, this time with the

predictions of the variable capital utilization model with just the

investment wedge. Comparing this to Figure 3, we see that with variable

capital utilization, the investment wedge drives output the wrong way to an

even greater extent than in the benchmark economy.



In Figure 12, we compare the contributions of the sum of the efficiency and

labor wedges for the two specifications of capital utilization (fixed and

variable) during the Great Depression period. The figure shows that these

contributions are quite similar. While remeasuring the efficiency wedge

changes the relative contributions of the two wedges, it clearly has little

effect on their combined contribution.



Overall, then, taking account of variable capital utilization strengthens

our finding that in the Great Depression period, the efficiency and labor

wedges play a primary role and investment wedges do not.



\bigskip



\noindent b. \emph{Different Labor Supply Elasticities}



Now we consider the effects on our results of changing the elasticity of

labor supply. We assume in our benchmark model that preferences are

logarithmic in both consumption and leisure. Consider now an alternative

prototype economy with a different elasticity of labor supply. We show that

a result analogous to that in Proposition 3 holds: allowing for different

labor supply elasticities changes the size of the measured labor wedge but

not that of the measured investment wedge. Therefore, if the contribution of

the investment wedge is zero in the benchmark prototype economy, it is also

zero in an economy with a different labor supply elasticity.



To see that, consider an alternative prototype economy which is identical to

a deterministic version of our benchmark model except that now the utility

function is given by $U(c)+V_{2}(1-l)$. Denote the utility function in our

benchmark prototype economy (economy 1) by $U(c)+V_{1}(1-l).$ Clearly, by

varying the function $V_{2},$ we can generate a wide range of alternative

labor supply elasticities.



Let the sequence of wedges and the equilibrium outcomes in the two economies

be $(A_{it},\tau _{lit},\tau _{xit})$ and $(y_{it},c_{it},l_{it},x_{it})$

for $i=1,2.$ We then have this proposition:



\bigskip



\textsc{Proposition }4: \emph{If the sequence of wedges for the alternative

prototype economy, economy }2,\emph{\ is given by } 

\[

1-\tau _{l2t}=(1-\tau _{l1t})\frac{V_{2}^{\prime }(1-l_{1t})}{V_{1}^{\prime

}(1-l_{1t})}, 

\]%

\emph{and if} $A_{2t}=A_{1t}$ \emph{and} $\tau _{x2t}=\tau _{x1t},$ \emph{%

then the equilibrium outcomes for the two economies coincide. }\bigskip



\textsc{Proof:} We prove this proposition by showing that the equilibrium

conditions of economy $2$ are satisfied at the equilibrium outcomes of

economy $1.$ The first-order condition for labor input in economy $1$ is 

\[

-\frac{V_{1}^{\prime }(1-l_{1t})}{U^{\prime }(c_{1t})}=(1-\tau _{l1t})\frac{%

(1-\alpha )y_{1t}}{l_{1t}}. 

\]%

Using the definition of $\tau _{l2t},$ we have that 

\[

-\frac{V_{2}^{\prime }(1-l_{1t})}{U^{\prime }(c_{1t})}=(1-\tau _{l2t})\frac{%

(1-\alpha )y_{1t}}{l_{1t}}, 

\]%

so that the first-order condition for labor in economy 2 is satisfied. The

rest of the equations governing the equilibrium are unaffected. 

%TCIMACRO{\TeXButton{TeX field}{\hspace{4.5in}} }%

%BeginExpansion

\hspace{4.5in}

%EndExpansion

$Q.E.D.$



\bigskip



Note that here even if the labor wedge does not fluctuate in our benchmark

prototype economy, it typically will in the alternative prototype economy.

Note also that the investment wedges are the same in both economies. Thus,

if the investment wedge is constant in one economy, it is constant in the

other; and the contribution of the investment wedge to fluctuations is zero

in both economies. Extending this proposition to a stochastic environment is

immediate.



\bigskip



\noindent c. \emph{Investment Adjustment Costs}



Now we consider a third alternative prototype economy, this one with

investment adjustment costs. These costs can be interpreted as standing in

for one of two features of detailed economies. One is that the detailed

economies have adjustment costs in converting output into installed capital.

Another interpretation is that the detailed model does not have adjustment

costs, but that financial frictions manifest themselves as adjustment costs

in the alternative prototype economy (as in Bernanke and Gertler (1989) and

Carlstrom and Fuerst (1997)).



In this alternative prototype economy, the only difference from the

benchmark prototype economy is that the capital accumulation law is no

longer (\ref{1}), but rather is%

\begin{equation}

(1+\gamma _{n})k_{t+1}=(1-\delta )k_{t}+x_{t}-\phi \left( \frac{x_{t}}{k_{t}}%

\right) k_{t},  \label{capital adjust}

\end{equation}%

where $\phi $ represents the per unit cost of adjusting the capital stock.

In the macroeconomic literature, a commonly used functional form for the

adjustment costs $\phi $ is 

\begin{equation}

\phi \left( \frac{x}{k}\right) =\frac{a}{2}\left( \frac{x}{k}-b\right) ^{2},

\label{adjustment}

\end{equation}%

where $b=\delta +\gamma +\gamma _{n}$ is the steady-state value of the

investment-capital ratio.



To set up the analog of Propositions 3 and 4, let the wedges in the

benchmark prototype economy, economy 1, and the alternative prototype

economy, economy 2, be $(A_{it},\tau _{lit},\tau _{xit})$ and $%

(y_{it},c_{it},l_{it},x_{it})$ for $i=1,2.$ For simplicity, let $\tau _{x2t}$

and $g_{2t}$ be identically zero. The proof of the following proposition is

immediate.



\bigskip



\textsc{Proposition }5: \emph{If the sequence of wedges for the alternative

prototype economy, economy }2,\emph{\ is given by }$A_{2t}=A_{1t},\tau

_{l2t}=\tau _{l1t},$ $\tau _{x1t}$ implicitly defined by\emph{\ } 

\begin{equation}

\frac{A_{1t+1}F_{kt+1}+(1-\delta )(1+\tau _{x1t+1})}{1+\tau _{x1t}}=\frac{%

A_{2t+1}F_{kt+1}+(1-\delta +\phi _{t}-x_{2t+1}\phi _{t+1}^{\prime

}/k_{2t+1})(1-\phi _{t+1}^{\prime })^{-1}}{(1-\phi _{t}^{\prime })^{-1}},

\label{phi}

\end{equation}%

\emph{and }$g_{1t}=\phi _{t}k_{2t},$ \emph{then the equilibrium outcomes for

the two economies coincide}$.$ \bigskip



In understanding the investment wedge $\tau _{x1t}$, note that if adjustment

costs are given by (\ref{adjustment}), then the term $\phi _{t}-x_{t+1}\phi

_{t+1}^{\prime }/k_{t+1}$ in (\ref{phi}) equals $a[(x/k)^{2}-b^{2}]/2$,

which is an order of magnitude smaller than $\phi _{t}^{\prime }=a[x/k-b].$

Setting this term to zero gives the approximation 

\begin{equation}

1+\tau _{x1t}=\frac{1}{1-\phi _{t}^{\prime }}.  \label{approx}

\end{equation}%

>From (\ref{approx}) and the convexity of $\phi $, we see that $\tau _{x1}$

is increasing in $x/k$ and is zero when $x/k$ is at its steady-state value.

Hence, in recessions, when $x/k$ is relatively low, $\tau _{x1}$ is

negative; and in booms, when $x/k$ is relatively high, $\tau _{x1}$ is

positive. In this sense, even when the alternative prototype economy has no

investment wedges, $\tau _{x1}$ will be countercyclical in the benchmark

prototype economy, so investment distortions will be smaller in recessions.



Note, more generally, that when investment wedges in the alternative

prototype economy are nonzero, the analog of (\ref{approx}) is 

\begin{equation}

1+\tau _{x1t}=\frac{1+\tau _{x2t}}{1-\phi _{t}^{\prime }}.  \label{approx2}

\end{equation}



We can also use the equivalence map in (\ref{approx2}) in the reverse

direction. Imagine that the data are generated from a detailed economy with

no adjustment costs and no frictions. The benchmark prototype economy will

have $\tau _{x1t}=0.$ Suppose a researcher considers an alternative

prototype economy with adjustment costs. This researcher will find that $%

1+\tau _{x2t}=1-\phi _{t}^{\prime }$, so that the investment wedge will be

procyclical even though the detailed economy has no frictions.



More generally, a researcher who incorrectly specifies too high a level of

adjustment costs in the alternative prototype economy will infer that

investment wedges play a much larger role than they actually do. To get some

intuition for this result, consider two alternative prototype economies $A$

and $B$, both of which have investment wedges and adjustment costs. The

analog of our approximation (\ref{approx2}) is $(1+\tau _{xAt})/(1-\phi

_{At}^{\prime })=(1+\tau _{xBt})/(1-\phi _{Bt}^{\prime })$, where $\phi

_{it} $ denotes adjustment costs in economy $i=A,B.$ Straightforward algebra

establishes that specifying too high a level of adjustment costs makes

investment wedges seem worse in recessions than they actually are.



Now we consider a prototype economy with adjustment costs and ask whether

our substantive conclusion that the investment wedge plays at most a

tertiary role in either the Great Depression or postwar business cycles is

affected. We follow Bernanke, Gertler, and Gilchrist (1999) in how we choose

the value for the parameter $a.$ Bernanke, Gertler, and Gilchrist (BGG)

choose this parameter so that the elasticity, $\eta ,$ of the price of

capital with respect to the investment-capital ratio is $.25.$ In this

setup, the price of capital $q=1/(1-\phi ^{\prime }),$ so that, evaluated at

the steady state, $\eta =a(\delta +\gamma +\gamma _{n}).$ Given our other

parameters, $a=3.22.$ In Figures 13 and 14, this parameterization is labeled

as the model with \emph{BGG Costs. }Bernanke, Gertler, and Gilchrist also

argue that a reasonable range for the elasticity $\eta $ is between $0$ and $%

.5$, that values much outside this range imply implausibly high adjustment

costs. We consider an extreme case in which $\eta =1$, so that $a=12.88$,

roughly four times the BGG level. In the figures, this parameterization is

labeled as the model with \emph{Extreme Costs. }For each setting of the

parameter $a$, we reestimate the stochastic process for the state. (For the

parameter values of the stochastic process, see the technical appendix.)



Comparing the investment wedges in Figure 1 and panel A of Figure 13, we see

that introducing investment adjustment costs leads the investment wedge to

worsen rather than improve in the early part of the Great Depression. In

panel A of Figure 14 we see that this worsening produces a decline in output

from 1929 to 1933. With the BGG adjustment costs, however, the decline is

tiny $(2\%$ from 1929 to 1933). Even with the extreme adjustment costs, the

decline is only $6.5\%$ and, hence, accounts for only about one-sixth of the

overall decline in output.\footnote{%

In general, in our Great Depression experiments, linear methods perform

poorly compared to our nonlinear method. With extreme adjustment costs, for

example, linear methods produce large errors, on the order of $100\%$.}

Moreover, with these extreme adjustment costs, the consumption anomaly

associated with investment wedges is acute. For example, from 1929 to 1932,

relative to trend, consumption in the data falls about $18\%$, while in the

model it actually rises by more than $5\%.$



We are skeptical that detailed models with extreme adjustment costs are

worth exploring. With extreme costs, given the observed levels of investment

and the capital stock in the data, (\ref{adjustment}) implies that the

resources lost due to adjustment costs as a fraction of output are nearly $%

7\%$ in 1933. We share Bernanke, Gertler, and Gilchrist's (1999) concerns

that costs of this magnitude are implausibly large. From 1929 to 1933,

investment falls sharply, but the adjustment costs implied by (\ref%

{adjustment}) rise sharply. Why would firms incur adjustment costs simply by

investing at positive rates below their steady-state value? The idea that

managers incurred huge adjustment costs simply because they were watching

their machines depreciate seems farfetched. Furthermore, the idea that from

1929 to 1933 investment fell sharply but adjustment costs rose sharply is

inconsistent with interpreting these costs as arising from monitoring costs

in an economy with financial frictions. Indeed, in such an economy, as

investment activity falls, so do monitoring costs.



We do similar experiments for the 1982 recession period, with similar

results. In panel B of Figure 13, we see that with BGG costs the investment

wedge fluctuates little through the end of 1982 and improves thereafter.

With extreme costs the investment wedges worsens until 1983 and improves

thereafter. Panel B of Figure 14 shows that in the model with costs at the

BGG level, the investment wedge plays essentially no role in output

fluctuations. With costs at four times the BGG level, the panel shows that

the investment wedge plays a bit larger but still modest role. Recall that,

relative to trend, output in the data falls almost 10\% from its peak to its

trough. With extreme adjustment costs, output due to the investment wedge

falls about 2.2\%, so that the investment wedge appears to account for

roughly one-fifth of the fall in output.



In our judgment, the investment wedge actually accounts for much less than

one-fifth of the 1982 recession fall in output because we think extreme

adjustment costs are implausible. Even aside from the Great Depression

issues, we find models with extreme adjustment costs difficult to reconcile

with data on plant level investment decisions. An extensive literature has

documented that investment at the plant level has large spikes. Doms and

Dunne (1994), for example, examine plant level data for 33,000 plants over

17 years. They show that much of the growth in the capital stock of a plant

is concentrated in a short period of time. In the year of highest

investment, the capital stock grows 45\%, while in the years immediately

before and after that, it grows less than 10\%. Such behavior is clearly

inconsistent with extreme adjustment costs.



Furthermore, the use of adjustment costs in macroeconomic analysis is

controversial. Kydland and Prescott (1982), for example, have argued that

models with adjustment costs like those in equation (\ref{adjustment}) are

inconsistent with the data. Such models imply a static relationship between

the investment-capital ratio and the relative price of investment goods to

output (of the form $q=1/[1-\phi ^{\prime }(x/k)]$). This means that the

elasticity of the investment-capital ratio with respect to the relative

price is the same in the short run and the long run. Kydland and Prescott

argue that this is not consistent with the data, where short-run

elasticities are much smaller than long-run elasticities.



Finally, the finding that investment wedges play a modest role when

adjustment costs are extreme is easy to reconcile with the view that

financial frictions which manifest themselves primarily as investment wedges

play a small role over the business cycle. This reconciliation uses the

insight discussed above that a modeler who incorrectly specifies too high a

level of adjustment costs will incorrectly find too large a role for

investment wedges.



\bigskip



\begin{center}

5. \textsc{Contrasting Our Decomposition With Traditional Decompositions}

\end{center}



Our decomposition of business cycle fluctuations is intended to isolate the

partial effects of each of the wedges on equilibrium outcomes, and in this

sense, it is different from traditional decompositions. Those decompositions

attempt to isolate the effects of (so-called) primitive shocks on

equilibrium outcomes; ours does not. Isolating the effects of primitive

shocks requires specifying a detailed model. Since our procedure precedes

the specification of a detailed model, it obviously cannot be used to

isolate the effects of primitive shocks. In order to clarify the distinction

between our decomposition and a traditional one, here we describe a

traditional decomposition and explain why we prefer ours.



The traditional decomposition attempts to isolate the effects of primitive

shocks by \textquotedblleft naming the innovations.\textquotedblright\

Recall that in our stochastic process for the four wedges, (\ref{s process}%

), the innovations $\varepsilon _{t+1}$ are allowed to be contemporaneously

correlated with the covariance matrix $V.$ Under the traditional

decomposition, the primitive shocks, say, $\eta _{t+1}$, are assumed to be

mean zero, to be contemporaneously uncorrelated with $E\eta _{t+1}\eta

_{t+1}^{\prime }=I$, and to lead to the same stochastic process for the

wedges. Identifying these primitive shocks requires specifying a matrix $R$

so that $R\eta _{t+1}=\varepsilon _{t+1}$ and $RR^{\prime }=V.$ Names are

then made up for these shocks, including \emph{money} shocks, \emph{demand}

shocks, \emph{technology} shocks, and so on.



In the traditional method, then, given any sequence of realized wedges $%

s_{t} $ and the specification of the matrix $R$, the associated realized

values of $\eta _{t}=(\eta _{1t},\eta _{2t},\eta _{3t},\eta _{4t})^{\prime }$

are computed. The movements in, say, output, are then decomposed into the

movements due to each one of these primitive shocks as follows. Let $%

s_{t}(\eta _{1})=$ {\large (}$\log A_{t}(\eta _{1}),\tau _{lt}(\eta _{1})$, $%

\tau _{xt}(\eta _{1}),\log g_{t}(\eta _{1})${\large )} denote the realized

values of the four wedges when the primitive shock sequence $\eta _{t}=(\eta

_{1t},0,0,0)$ is fed into 

\[

s_{t+1}=P_{0}+Ps_{t}+R\eta _{t+1}. 

\]%

(Note that when the covariance matrix $V$ is not diagonal, movements in a

primitive shock $\eta _{1t}$ will lead to movements in more than one wedge.)

The predicted value of, say, output due to $\eta _{1}$ is then computed from

the decision rules of the model according to $y_{t}(\eta _{1})=y${\large (}$%

k_{t}(\eta _{1}),s_{t}(\eta _{1})${\large )}, where $k_{t}(\eta _{1})$ is

computed recursively using the decision rule for investment, the initial

capital stock, and the capital accumulation law. The values $s_{t}(\eta

_{2}) $, $y_{t}(\eta _{2})$, and the rest are computed in a similar way.



Notice that in this traditional decomposition method, the realized value of

each wedge is simply the sum of the parts due to $\eta _{1},\eta _{2},$ and

so on, in that 

\[

\log A_{t}=\sum_{i}\log A_{t}(\eta _{i}),~\tau _{lt}=\sum_{i}\tau _{lt}(\eta

_{i}),~\tau _{xt}=\sum_{i}\tau _{xt}(\eta _{i}),~\log g_{t}=\sum_{i}\log

g_{t}(\eta _{i}). 

\]%

In this sense, the traditional decomposition attempts to decompose each of

the four wedges into four component parts, each of which is due to a

primitive shock.



Our decomposition is purposefully less ambitious. It includes only the

effect of the movements of each total wedge, not any of its subparts. The

advantage of our decomposition is that it is invariant to $R$, so that we do

not need to make identifying assumptions implicit in the specification of

the matrix $R,$ nor do we need to make up names for the shocks. The

invariance makes our method valuable. The problem with the traditional

approach is that finding identifying assumptions that apply to a broad class

of detailed models is very difficult. Hence, this approach is not useful in

pointing researchers toward classes of promising models. In a sense, the

traditional approach puts the cart before the horse: the decomposition

requires specifying the matrix $R$, but doing that requires a detailed

model. Our decomposition is useful precisely because it does not need to

make up identifying assumptions to specify the matrix $R.$ Our equivalence

results have identifying assumptions built into them, so that results from

the benchmark prototype economy can be used to uncover promising classes of

detailed models.



\bigskip



\begin{center}

6. \textsc{Reviewing the Related Literature}

\end{center}



Our work here is related to the existing literature in terms of methodology

and the interpretation of the wedges.



\bigskip \newpage



\begin{center}

6.1. \emph{Methodology}

\end{center}



Our basic method is to use restrictions from economic theory to back out

wedges from the data, formulate stochastic processes for these wedges, and

then put the wedges back into a quantitative general equilibrium model for

an accounting exercise. This basic idea is at the heart of an enormous

amount of work in the real business cycle theory literature. Prescott

(1986), for example, asks what fraction of the variance of output can

plausibly be attributed to productivity shocks, which we have referred to as

the \emph{efficiency wedge}. Subsequent studies have expanded this general

equilibrium accounting exercise to include a wide variety of other shocks.

(See, for example, the studies in Cooley's 1995 volume.)



An important difference between our method and others is that we back out

the labor wedge and the investment wedge from the combined consumer and firm

first-order conditions, while most of the recent business cycle literature

uses direct measures of labor and investment shocks. Perhaps the most

closely related precursor of our method is McGrattan's (1991); she uses the

equilibrium of her model to infer the implicit wedges. Ingram, Kocherlakota,

and Savin (1997) advocate a similar approach.



\bigskip



\begin{center}

6.2. \emph{Wedge Interpretations}

\end{center}



The idea that taxes of various kinds distort the relation between various

marginal conditions is the cornerstone of public finance. Taxes are not the

only well-known distortions; monopoly power by unions or firms is also

commonly thought to produce a labor wedge. And the idea that a labor wedge

is produced by sticky wages or sticky prices is the cornerstone of the new

Keynesian approach to business cycles; see, for example, the survey by

Rotemberg and Woodford (1999). One contribution of our work here is to show

the precise mapping between the wedges and general equilibrium models with

frictions.



Many studies have plotted one or more of the four wedges. The efficiency

wedge has been extensively studied. (See, for example, Kehoe and Prescott

(2002).) The labor wedge has also been studied. For example, Parkin (1988),

Hall (1997), and Gal\'{\i}, Gertler, and L\'{o}pez-Salido (forthcoming) all

graph and interpret the labor wedge for the postwar data. Parkin discusses

how monetary shocks might drive this wedge, and Hall discusses how search

frictions might drive it. Gal\'{\i}, Gertler, and L\'{o}pez-Salido discuss a

variety of interpretations of the labor wedge, as do Rotemberg and Woodford

(1991, 1999). Mulligan (2002a, b) plots the labor wedge for the United

States for much of the 20th century, including the Great Depression period.

He interprets movements in this wedge as arising from changes in labor

market institutions and regulation, including features we discuss here. Cole

and Ohanian (2002) plot the labor wedge for the Great Depression and offer

interpretations similar to ours. The investment wedge has been investigated

by McGrattan (1991), Braun (1994), Carlstrom and Fuerst (1997), and Cooper

and Ejarque (2000).



\bigskip



\begin{center}

7. \textsc{Conclusions and Extensions}

\end{center}



This study is aimed at applied theorists who are interested in building

detailed, quantitative models of economic fluctuations. Once such theorists

have chosen the primitive sources of shocks to economic activity, they need

to choose the mechanisms through which the shocks lead to business cycle

fluctuations. We have shown that these mechanisms can be summarized by their

effects on four wedges in the standard growth model. Our business cycle

accounting method can be used to judge which mechanisms are promising and

which are not, thus helping theorists narrow their options. We view our

method as an alternative to the use of structural vector autoregressions

(VARs), which has also been advocated as a way to identify promising

mechanisms. (Elsewhere, in Chari, Kehoe, and McGrattan (2005), we argue that

structural VARs have deficiencies that limit their usefulness.)



Here we have demonstrated how our method works by applying it to two

historical episodes, the Great Depression and the 1982 U.S. recession. We

have found that efficiency and labor wedges, in combination, account for

essentially all of the decline and recovery in these business cycles;

investment wedges play, at best, a tertiary role. These results hold in

summary statistics of the entire postwar period and in alternative

specifications of the growth model. We have also found that when we maximize

the contribution of the investment wedge, the models display a consumption

anomaly: they tend to produce much smaller declines in consumption during

downturns than occur in the data. These findings together imply that

existing models of financial frictions in which the distortions primarily

manifest themselves as investment wedges can account, at best, for only a

small fraction of the fluctuations in the Great Depression or more typical

U.S. downturns. This finding is our primary substantive contribution.



We have seen that if adjustment costs are extreme, investment wedges can

play a larger but still modest role over the business cycle. A combination

of microeconomic and macroeconomic observations argue against extreme

adjustment costs, and we view them as a theoretical curiosity of no applied

interest.



A useful extension of our work here would be to decompose our wedges into a

portion that comes from explicit taxes imposed by governments and a portion

that comes from frictions in detailed models. This decomposition would be

particularly useful when explicit taxes vary significantly over the business

cycle.



Researchers have argued---quite beyond the prominent model of Bernanke and

Gertler (1989)---that frictions in financial markets are important for

business cycle fluctuations. (See Bernanke (1983) and the motivation in

Bernanke and Gertler (1989).) We stress that our findings do not contradict

this idea. Indeed, we have shown that a detailed economy with

input-financing frictions is equivalent to a prototype economy with

efficiency wedges. In this sense, while existing models of financial

frictions are not promising, new models in which financial frictions show up

as efficiency and labor wedges are.



Our results suggest that future theoretical work should focus on developing

models which lead to fluctuations in efficiency and labor wedges. Many

existing models produce fluctuations in labor wedges. The challenging task

is to develop detailed models in which primitive shocks lead to fluctuations

in efficiency wedges as well.



\bigskip



\noindent \emph{University of Minnesota and Research Department, Federal

Reserve Bank of Minneapolis, 90 Hennepin Avenue, Minneapolis, MN 55401,

U.S.A.; chari@res.mpls.frb.fed.us}



\begin{center}

and

\end{center}



\emph{Research Department, Federal Reserve Bank of Minneapolis, 90 Hennepin

Avenue, Minneapolis, MN 55401, U.S.A., and University of Minnesota;

pkehoe@res.mpls.frb.fed.us}



\begin{center}

and

\end{center}



\emph{Research Department, Federal Reserve Bank of Minneapolis, 90 Hennepin

Avenue, Minneapolis, MN 55401, U.S.A., and University of Minnesota;

erm@mcgrattan.mpls.frb.fed.us}



\bigskip \newpage



%TCIMACRO{%

%\TeXButton{TeX field}{\renewcommand{\baselinestretch}{1.0}\normalsize}}%

%BeginExpansion

\renewcommand{\baselinestretch}{1.0}\normalsize%

%EndExpansion



\begin{center}

\textbf{APPENDIX: The Mapping for Two Other Wedges}

\end{center}



Here we demonstrate the mapping from two other detailed economies with

frictions to two prototype economies with wedges. In the preceding text, we

have described the mapping for efficiency and labor wedges. In this

appendix, we describe it for investment and government consumption wedges.



\bigskip



\begin{center}

\textbf{A. Investment Wedges Due to Financial Frictions}

\end{center}



We start with the mapping of financial frictions to investment wedges.



In Chari, Kehoe, and McGrattan (2004), we show the equivalence between the

Carlstrom and Fuerst (1997) model and a prototype economy. Here we focus on

the financial frictions in the Bernanke, Gertler, and Gilchrist (1999) model

and abstract from the monetary features of that model. Bernanke, Gertler,

and Gilchrist begin by deriving the optimal financial contracts between

risk-neutral entrepreneurs and financial intermediaries in an environment

with no aggregate uncertainty. These contracts resemble debt contracts (with

default). Bernanke, Gertler, and Gilchrist try to extend their derivation of

optimal contracts to an economy with aggregate uncertainty. The contracts

they consider are not optimal given the environment, because these contracts

do not allow risk sharing between risk-averse consumers and risk-neutral

entrepreneurs. Here we solve for optimal debt-like contracts that

arbitrarily rule out such risk sharing, but even so, our first-order

conditions differ from those of Bernanke, Gertler, and Gilchrist. (One

reason for this difference is that our break-even constraint for the

financial intermediary and our law of motion for entrepreneurial net worth

differ from theirs. See equations (\ref{breakeven}) and (\ref{n})\ below.)



\bigskip



\noindent a. \emph{A Detailed Economy With Financial Frictions}



The Bernanke-Gertler-Gilchrist model has a continuum of risk-neutral

entrepreneurs of mass $L_{e}$, a continuum of consumers of mass 1, and a

representative firm. Output $y(s^{t})$ is produced according to 

\begin{equation}

y(s^{t})=A(s^{t})k^{\alpha }\left[ l^{\Omega }L_{e}^{1-\Omega }\right]

^{1-\alpha },  \label{bggprodn}

\end{equation}%

where $A(s^{t}),k,$ and $l$ are the technology shock, the capital stock, and

the labor supplied by consumers. The stochastic process for the technology

shock is given by $\log A(s^{t+1})=(1-\rho _{A})\log A+\rho _{A}\log

A(s^{t})+\varepsilon _{A}(s^{t+1}).$ Each entrepreneur supplies one unit of

labor inelastically. The representative firm's maximization problem is to

choose $l,L_{e},k$ to maximize profits $A(s^{t})k^{\alpha }\left[ l^{\Omega

}L_{e}^{1-\Omega }\right] ^{1-\alpha }-w(s^{t})l-w_{e}(s^{t})L_{e}-r(s^{t})k$%

, where $w(s^{t}),$ $w_{e}(s^{t}),$ and $r(s^{t})$ denote the wage rate for

consumers, the wage rate for entrepreneurs, and the rental rate on capital.



New capital goods can be produced only by entrepreneurs. Each entrepreneur

owns a technology that transforms output and old capital at the end of any

period into capital goods at the beginning of the following period. In each

period $t,$ each entrepreneur receives an idiosyncratic shock $\omega $

drawn from a distribution $F(\omega )$ with expected value 1. This shock is

i.i.d. across entrepreneurs and time. The realization of $\omega $ is

private information to the entrepreneur. An entrepreneur who buys $%

k(s^{t-1}) $ units of goods in period $t-1$ produces $\omega k(s^{t-1})$

units of capital at the beginning of period $t.$ These capital goods are

sold at price $R_{k}(s^{t})=r(s^{t})+(1-\delta )$, where $\delta $ is the

rate of depreciation of the capital goods.



Entrepreneurs finance the production of new capital goods partly with their

own net worth, $n(s^{t-1}),$ and partly with loans from financial

intermediaries. These intermediaries offer contracts with the following

cutoff form: in all idiosyncratic states at $t$ in which $\omega \geq \bar{%

\omega}(s^{t})$, the entrepreneur pays $\bar{\omega}%

(s^{t})R_{k}(s^{t})k(s^{t-1})$ and keeps $[\omega -\bar{\omega}%

(s^{t})]R_{k}(s^{t})k(s^{t-1}).$ In all other states, the entrepreneur

receives nothing while the financial intermediary receives $\omega

R_{k}(s^{t})k(s^{t-1})$ net of monitoring costs $\mu \omega

R_{k}(s^{t})k(s^{t-1}).$ We assume that financial intermediaries make zero

profits in equilibrium.



Given an entrepreneur's net worth $n(s^{t-1}),$ the contracting problem for

a representative entrepreneur is to choose the cutoff $\bar{\omega}(s^{t})$

for each state and the amount of goods invested $k(s^{t-1})$ to maximize the

expected utility of the entrepreneurs, 

\[

\sum \pi (s^{t}|s^{t-1})\text{{\large \{[}}(1-\Gamma \text{{\large (}}\bar{%

\omega}(s^{t})\text{{\large )]}}R_{k}(s^{t})k(s^{t-1})~dF(\omega )\text{%

{\large \}}}, 

\]

subject to a break-even constraint for the intermediary, 

\begin{equation}

\sum (s^{t}|s^{t-1})\text{{\large [}}\Gamma \text{{\large (}}\bar{\omega}%

(s^{t})\text{{\large )}}-\mu G\text{{\large (}}\bar{\omega}(s^{t})\text{%

{\large )]}}R_{k}(s^{t})k(s^{t-1})=k(s^{t-1})-n(s^{t-1}),  \label{breakeven}

\end{equation}

where $\Gamma (\bar{\omega})=\int_{0}^{\bar{\omega}}\omega f(\omega

)~d\omega +\bar{\omega}\int_{\bar{\omega}}^{\infty }f(\omega )~d\omega $ and 

$G(\bar{\omega})=\int_{0}^{\bar{\omega}}\omega f(\omega )~d\omega .$ Here $%

q(s^{t}|s^{t-1})$ denotes the price of consumption goods in state $s^{t}$ in

units of consumption goods in state $s^{t-1}.$



This formulation implies an aggregation result: The capital demand of the

entrepreneurs is linear in their net worth, so that total demand for goods

from them depends only on their total net worth. The solution to the

contracting problem is characterized by the first-order conditions with

respect to $k(s^{t-1})$ and $\bar{\omega}(s^{t})$, which can summarized by 

\begin{equation}

\frac{\sum \pi (s^{t}|s^{t-1})\text{{\large [}}1-\Gamma \text{{\large (}}%

\bar{\omega}(s^{t})\text{{\large )]}}R_{k}(s^{t})}{\sum \pi

(s^{t}|s^{t-1})\Gamma ^{\prime }\text{{\large (}}\bar{\omega}(s^{t})\text{%

{\large )}}R_{k}(s^{t})}=\frac{\sum q(s^{t}|s^{t-1})\text{{\large [}}\Gamma 

\text{{\large (}}\bar{\omega}(s^{t})\text{{\large )}}-\mu G\text{{\large (}}%

\bar{\omega}(s^{t})\text{{\large )]}}R_{k}(s^{t})-1}{\sum q(s^{t}|s^{t-1})%

\text{{\large [}}\Gamma ^{\prime }\text{{\large (}}\bar{\omega}(s^{t})\text{%

{\large )}}-\mu G^{\prime }\text{{\large (}}\bar{\omega}(s^{t})\text{{\large %

)]}}R_{k}(s^{t})}  \label{OmegaBar}

\end{equation}

and the break-even constraint (\ref{breakeven}).



In each period, a fraction $\gamma $ of existing entrepreneurs dies and is

replaced by a fraction $\gamma $ of newborn entrepreneurs. At the beginning

of each period, entrepreneurs learn whether they will die this period. We

assume that $\gamma $ is sufficiently small and the technology for producing

investment goods is sufficiently productive so that entrepreneurs consume

their net worth only when they are about to die. All entrepreneurs who are

going to die consume their entire net worth and do not supply labor in the

current period. (The newborn replacements do instead.) Those entrepreneurs

who do not die save their wage income plus their income from producing

capital goods. The aggregate income of the entrepreneurs is, then, 

\[

\int \left[ \omega -\bar{\omega}(s^{t})\right] R_{k}(s^{t})k(s^{t-1})~dF(%

\omega )+w_{e}(s^{t})L_{e}, 

\]

where $k(s^{t-1})$ is the aggregate capital stock. The law of motion for

aggregate net worth is given by



\begin{equation}

n(s^{t})=\gamma \int \left[ \omega -\bar{\omega}(s^{t})\right]

R_{k}(s^{t})k(s^{t-1})~dF(\omega )+w_{e}(s^{t})L_{e}.  \label{n}

\end{equation}

Total consumption by entrepreneurs in any period is 

\begin{equation}

c_{e}(s^{t})=(1-\gamma )\int \left[ \omega -\bar{\omega}(s^{t})\right]

R_{k}(s^{t})k(s^{t-1})~dF(\omega ).  \label{ce}

\end{equation}



Consumers maximize utility given by 

\begin{equation}

\sum_{t=0}^{\infty }\sum_{s^{t}}\,\beta ^{t}\pi _{t}(s^{t})\,U\text{{\large (%

}}c(s^{t}),l(s^{t})\text{{\large )}},  \label{utility}

\end{equation}%

where $c(s^{t})$ denotes consumption in state $s^{t}$ subject to 

\begin{equation}

c(s^{t})+\sum_{s_{t+1}}q(s^{t+1}|s^{t})b(s^{t+1})\leq

w(s^{t})l(s^{t})+b(s^{t})+T(s^{t}),  \nonumber

\end{equation}%

for $t=0,1,\ldots ,$ and borrowing constraints, $b(s^{t+1})\geq $ $\bar{b},$

for some large negative number $\bar{b}.$ Here $T(s^{t})$ is lump-sum

transfers. The initial condition $b(s^{0})$ is given. Each of the bonds $%

b(s^{t+1})$ is a claim to one unit of consumption in state $s^{t+1}$ and

costs $q(s^{t+1}|s^{t})$ dollars in state $s^{t}.$ The first-order

conditions for the consumer can be written as 

\begin{eqnarray}

-\frac{U_{l}(s^{t})}{U_{c}(s^{t})} &=&w(s^{t})\text{ and}

\label{Euler for money} \\

q(s^{t+1}|s^{t}) &=&\beta \pi (s^{t+1}|s^{t})\frac{U_{c}(s^{t+1})}{%

U_{c}(s^{t})}.  \label{interest rate}

\end{eqnarray}



The market-clearing condition for final goods is then 

\[

c(s^{t})+c_{e}(s^{t})+\mu G\text{{\large (}}\bar{\omega}(s^{t})\text{{\large %

)}}R_{k}(s^{t})k(s^{t-1})+k(s^{t})-(1-\delta )k(s^{t-1})=y(s^{t}). 

\]



\bigskip



\noindent b. \emph{The Associated Prototype Economy With Investment Wedges}



Now consider a prototype economy which is the same as our benchmark

prototype economy, except for the following changes. We assume that the

production function is as in (\ref{bggprodn}), with $L_{e}$ interpreted as a

fixed input rented from the government, and we assume that consumers are

taxed on capital income but not on investment or labor income. With capital

income taxes, the consumers' budget constraint is given by 

\[

c(s^{t})+k(s^{t})-(1-\delta )k(s^{t-1})=w(s^{t})l(s^{t})+[1-\tau

_{k}(s^{t})]r(s^{t})k(s^{t-1})+\delta \tau _{k}(s^{t})k(s^{t-1})+T(s^{t}), 

\]%

where $\tau _{k}(s^{t})$ is the tax rate on capital income. Here $\tau

_{k}(s^{t})$ plays the role of an investment wedge. The resource constraint

for the prototype economy is as in the benchmark prototype economy. Let

government consumption in the prototype economy be given by 

\begin{equation}

g(s^{t})=c_{e}^{\ast }(s^{t})+\mu G\text{{\large (}}\bar{\omega}^{\ast

}(s^{t})\text{{\large )}}R_{k}^{\ast }(s^{t})k^{\ast }(s^{t-1}),

\label{gbgg}

\end{equation}%

where asterisks denote the allocations in the detailed economy with

investment frictions.



We now show how the tax on capital income in our prototype economy can be

constructed from the detailed economy. First, the break-even constraint in

the detailed economy can be rewritten as 

\[

\sum q^{\ast }(s^{t}|s^{t-1})\left\{ \text{{\large [}}\Gamma \text{{\large (}%

}\bar{\omega}^{\ast }(s^{t})\text{{\large )}}-\mu G\text{{\large (}}\bar{%

\omega}^{\ast }(s^{t})\text{{\large )]}}R_{k}^{\ast }(s^{t})+\frac{n^{\ast

}(s^{t-1})}{k^{\ast }(s^{t-1})}\frac{U_{c}^{\ast }(s^{t-1})}{\beta \pi

(s^{t}|s^{t-1})U_{c}^{\ast }(s^{t})}\right\} =1, 

\]

where $q^{\ast }(s^{t}|s^{t-1})=\beta \pi (s^{t}|s^{t-1})U_{c}^{\ast

}(s^{t})/U_{c}^{\ast }(s^{t-1}).$ The intertemporal Euler equation in the

prototype economy is



\[

\sum q(s^{t}|s^{t-1})\text{{\large \{[}}(1-\tau _{k}(s^{t})\text{{\large ]}}%

A(s^{t})F_{k}(s^{t})+(1-\delta )\text{{\large \}}}=1, 

\]%

where $q(s^{t}|s^{t-1})=\beta \pi

(s^{t}|s^{t-1})U_{c}(s^{t})/U_{c}(s^{t-1}). $ Let the tax rate on capital

income $\tau _{k}(s^{t})$ be such that 

\begin{eqnarray}

&&[1-\tau _{k}(s^{t})][A(s^{t})F_{k}^{\ast }(s^{t})-\delta ]+1

\label{bggeuler} \\

&=&\text{{\large [}}\Gamma \text{{\large (}}\bar{\omega}^{\ast }(s^{t})\text{%

{\large )}}-\mu G\text{{\large (}}\bar{\omega}^{\ast }(s^{t})\text{{\large )]%

}}R_{k}^{\ast }(s^{t})+\frac{n^{\ast }(s^{t-1})}{k^{\ast }(s^{t-1})}\frac{%

U_{c}^{\ast }(s^{t-1})}{\beta \pi (s^{t}|s^{t-1})U_{c}^{\ast }(s^{t})}. 

\nonumber

\end{eqnarray}%

Comparing first-order conditions for the two economies, we then have this

proposition:



\bigskip



\textsc{Proposition }6:\emph{\ Consider the prototype economy just

described, with government consumption given by }(\ref{gbgg})\emph{\ and

capital income taxes }$\emph{by}$\emph{\ }(\ref{bggeuler})\emph{. The

aggregate equilibrium allocations for this prototype economy coincide with

those of the detailed economy with financial frictions}$.$



\bigskip



\begin{center}

\textbf{B. Government Consumption Wedges Due to International Borrowing and

Lending}

\end{center}



Now we develop a detailed economy with international borrowing and lending

and show that net exports in that economy are equivalent to a government

consumption wedge in an associated prototype economy.



\bigskip



\noindent a. \emph{A Detailed Economy With International Borrowing and

Lending}



Consider a model of a world economy with $N$ countries and a single

homogenous good in each period. We use the same notation for uncertainty as

before.



The representative consumer in country $i$ has preferences 

\begin{equation}

\sum \beta ^{t}\pi _{t}(s^{t})U\text{{\large (}}c_{i}(s^{t}),l_{i}(s^{t})%

\text{{\large )}},  \label{preferences 2}

\end{equation}%

where $c_{i}(s^{t})$ and $l_{i}(s^{t})$ denote consumption and labor. The

consumer's budget constraint is 

\begin{equation}

c_{i}(s^{t})+b_{i}(s^{t})+k_{i}(s^{t})\leq F\text{{\large (}}%

k_{i}(s^{t-1}),l_{i}(s^{t})\text{{\large )}}+(1-\delta

)k_{i}(s^{t-1})+\sum_{s^{t+1}}q(s^{t+1}/s^{t})b_{i}(s^{t+1}),  \label{16}

\end{equation}%

where $b^{i}(s^{t+1})$ denotes the amount of state-contingent borrowing by

the consumer in country $i$ in period $t$, $q(s^{t+1}/s^{t})$ denotes the

corresponding state-contingent price, and $k_{i}(s^{t})$ denotes the capital

stock.



An \textit{equilibrium for this detailed economy} is a set of allocations 

{\large (}$c_{i}(s^{t}),k_{i}(s^{t}),l_{i}(s^{t}),$\linebreak $%

b_{i}(s^{t+1}) ${\large )} and prices $q(s^{t}/s^{t-1})${\large \ }such that

these allocations both solve the consumer's problem in each country $i$ and

satisfy the world resource constraint: 

\begin{equation}

\sum_{i=1}^{N}[c_{i}(s^{t})+k_{i}(s^{t})]\leq \sum_{i=1}^{N}\text{{\large [}}%

F\text{{\large (}}k_{i}(s^{t-1}),l_{i}(s^{t})\text{{\large )}}+(1-\delta

)k_{i}(s^{t-1})\text{{\large ].}}  \label{world resource}

\end{equation}%

Note that in this economy, the net exports of country $i$ are given by 

\[

F\text{{\large (}}k_{i}(s^{t-1}),l_{i}(s^{t})\text{{\large )}}%

-[k_{i}(s^{t})-(1-\delta )k_{i}(s^{t-1})]-c_{i}(s^{t}). 

\]



\bigskip



\noindent b. \emph{The Associated Prototype Economy With Government

Consumption Wedges}



Now consider a prototype economy of a single closed economy $i$ with an

exogenous stochastic variable, government consumption\textit{\ }$%

g_{i}(s^{t}),$ which we call the \textit{government consumption wedge}. In

this economy, consumers maximize (\ref{preferences 2}) subject to their

budget constraint 

\begin{equation}

c_{i}(s^{t})+k_{i}(s^{t})=w_{i}(s^{t})l_{i}(s^{t})+\left[ r_{i}(s^{t})+1-%

\delta \right] k_{i}(s^{t-1})+T_{i}(s^{t}),  \label{prototype budget}

\end{equation}%

where $w_{i}(s^{t}),$ $r_{i}(s^{t}),$ and $T_{i}(s^{t})$ are the wage rate,

the capital rental rate, and lump-sum transfers. In each state $s^{t},$

firms choose $k$ and $l$ to maximize $F(k,l)-r_{i}(s^{t})k-w_{i}(s^{t})l$.

The government's budget constraint is 

\begin{equation}

g_{i}(s^{t})+T_{i}(s^{t})=0.  \label{government budget}

\end{equation}%

The resource constraint for this economy is 

\begin{equation}

c_{i}(s^{t})+g_{i}(s^{t})+k_{i}(s^{t})=F\text{{\large (}}%

k_{i}(s^{t-1}),l_{i}(s^{t})\text{{\large )}}+(1-\delta )k_{i}(s^{t-1}).

\label{resource 2}

\end{equation}%

An \textit{equilibrium of the prototype economy} is, then, a set of

allocations {\large (}$c_{i}(s^{t}),k_{i}(s^{t}),l_{i}(s^{t}),g_{i}(s^{t}),%

\linebreak T_{i}(s^{t})${\large )} and prices {\large (}$%

w_{i}(s^{t}),r_{i}(s^{t})${\large )} such that these allocations are optimal

for consumers and firms and the resource constraint is satisfied.



The following proposition shows that the government consumption wedge in

this prototype economy consists of net exports in the original economy.



\bigskip



\textsc{Proposition }7:\textit{\ }\emph{Consider the equilibrium allocations}

{\large (}$c_{i}^{\ast }(s^{t}),k_{i}^{\ast }(s^{t}),l_{i}^{\ast

}(s^{t}),b_{i}^{\ast }(s^{t+1})${\large )} \emph{for country }$i$\emph{\ in

the detailed economy. Let the government consumption wedge be} 

\begin{equation}

g_{i}(s^{t})=F\text{{\large (}}k_{i}^{\ast }(s^{t-1}),l_{i}^{\ast }(s^{t})%

\text{{\large )}}-\left[ k_{i}^{\ast }(s^{t})-(1-\delta )k_{i}^{\ast

}(s^{t-1})\right] -c_{i}^{\ast }(s^{t}),  \label{g}

\end{equation}%

\emph{let the wage and capital rental rates be }$w_{i}(s^{t})=F_{li}^{\ast

}(s^{t})$ \emph{and} $r_{i}(s^{t})=F_{ki}^{\ast }(s^{t}),$ \emph{and let} $%

T_{i}(s^{t})$ \emph{be defined by }(\ref{government budget}). \emph{Then the

allocations} {\large (}$c_{i}^{\ast }(s^{t}),k_{i}^{\ast

}(s^{t}),l_{i}^{\ast }(s^{t}),g_{i}^{\ast }(s^{t}),T_{i}(s^{t})${\large )} 

\emph{and the prices }{\large (}$w_{i}(s^{t}),r_{i}(s^{t})${\large )}\emph{\

are an equilibrium for the prototype economy.} \bigskip



The proof follows by noting that the first-order conditions are the same in

the two economies and that, given the government consumption wedge (\ref{g}%

), the consumer's budget constraint (\ref{16}) in the detailed economy is

equivalent to the resource constraint (\ref{resource 2}) in the prototype

economy.



Note that for simplicity we have abstracted from government consumption in

the detailed economy. If that economy had government consumption as well,

then the government consumption wedge in the prototype economy would be the

sum of net exports and government consumption in the detailed

economy.\newpage



%TCIMACRO{%

%\TeXButton{TeX field}{\renewcommand{\baselinestretch}{1.5}\normalsize}}%

%BeginExpansion

\renewcommand{\baselinestretch}{1.5}\normalsize%

%EndExpansion



\begin{center}

\textbf{\large REFERENCES}

\end{center}



\textsc{Bernanke, B.} (1983): \textquotedblleft Nonmonetary Effects of the

Financial Crisis in the Propagation of the Great

Depression,\textquotedblright\ \emph{American Economic Review,} 73, 257--276.



\textsc{Bernanke, B., and M. Gertler} (1989): \textquotedblleft Agency

Costs, Net Worth, and Business Fluctuations,\textquotedblright\ \textit{%

American Economic Review}, 79, 14--31.



\textsc{Bernanke, B., M. Gertler, and S. Gilchrist} (1999):

\textquotedblleft The Financial Accelerator in a Quantitative Business Cycle

Framework,\textquotedblright\ in \textit{Handbook of Macroeconomics,} Volume

1C, ed. by J. Taylor and M. Woodford. Amsterdam: North-Holland.



\textsc{Bordo, M., C. Erceg, and C. Evans} (2000): \textquotedblleft Money,

Sticky Wages, and the Great Depression,\textquotedblright\ \textit{American

Economic Review,} 90, 1447--1463.



\textsc{Braun, R.} (1994): \textquotedblleft How Large Is the Optimal

Inflation Tax?\textquotedblright\ \textit{Journal of Monetary} \textit{%

Economics,} 34, 201--214.



\textsc{Carlstrom, C., and T. Fuerst} (1997): \textquotedblleft Agency

Costs, Net Worth, and Business Fluctuations: A Computable General

Equilibrium Analysis,\textquotedblright\ \textit{American Economic\ Review,}

87, 893--910.



\textsc{Chari, V., P. Kehoe, and E. McGrattan} (2002): \textquotedblleft

Accounting for the Great Depression,\textquotedblright\ \textit{American

Economic Review}, 92, 22--27.



%TCIMACRO{\TeXButton{rule}{\rule[-0.02in]{0.81in}{0.2pt}} }%

%BeginExpansion

\rule[-0.02in]{0.81in}{0.2pt}

%EndExpansion

(2004): \textquotedblleft Business Cycle Accounting,\textquotedblright\

Working Paper 10351, National Bureau of Economic Research.



%TCIMACRO{\TeXButton{rule}{\rule[-0.02in]{0.81in}{0.2pt}} }%

%BeginExpansion

\rule[-0.02in]{0.81in}{0.2pt}

%EndExpansion

(2005): \textquotedblleft A Critique of Structural VARs Using Business Cycle

Theory,\textquotedblright\ Research Department Staff Report 364, Federal

Reserve Bank of Minneapolis.



%TCIMACRO{\TeXButton{rule}{\rule[-0.02in]{0.81in}{0.2pt}} }%

%BeginExpansion

\rule[-0.02in]{0.81in}{0.2pt}

%EndExpansion

(2006): \textquotedblleft Appendices: Business Cycle

Accounting,\textquotedblright\ Research Department Staff Report 362, Federal

Reserve Bank of Minneapolis.



\textsc{Christiano, L., C. Gust, and J. Roldos} (2004): \textquotedblleft

Monetary Policy in a Financial Crisis,\textquotedblright\ \emph{Journal of

Economic Theory,} 119, 64--103.



\textsc{Chu, T.} (2001): \textquotedblleft Exit Barriers and Economic

Stagnation: The Macro-impact of Government Subsidies to

Industries,\textquotedblright\ Manuscript, East-West Center.



\textsc{Cole, H., and L. Ohanian} (1999): \textquotedblleft The Great

Depression in the United States from a Neoclassical

Perspective,\textquotedblright\ \textit{Federal Reserve Bank of Minneapolis

Quarterly Review}, 23(Winter), 2--24.



%TCIMACRO{\TeXButton{rule}{\rule[-0.02in]{0.81in}{0.2pt}} }%

%BeginExpansion

\rule[-0.02in]{0.81in}{0.2pt}

%EndExpansion

(2002): \textquotedblleft The U.S. and U.K. Great Depressions through the

Lens of Neoclassical Growth Theory,\textquotedblright\ \textit{American

Economic Review,} 92, 28--32.



%TCIMACRO{\TeXButton{rule}{\rule[-0.02in]{0.81in}{0.2pt}} }%

%BeginExpansion

\rule[-0.02in]{0.81in}{0.2pt}

%EndExpansion

(2004): \textquotedblleft New Deal Policies and the Persistence of the Great

Depression: A General Equilibrium Analysis,\textquotedblright\ \textit{%

Journal of Political Economy}, 112, 779--816.



\textsc{Cooley, T.,}\ ed. (1995): \textit{Frontiers of Business Cycle

Research}. Princeton: Princeton University Press.



\textsc{Cooper, R., and J. Ejarque} (2000): \textquotedblleft Financial

Intermediation and Aggregate Fluctuations: A Quantitative

Analysis,\textquotedblright\ \ \textit{Macroeconomic Dynamics}, 4(4),

423--447.



\textsc{Doms, M., and T. Dunne }(1994): \textquotedblleft Capital Adjustment

Patterns in Manufacturing Plants\textquotedblright\ CES 94-11, U.S. Bureau

of the Census.



\textsc{Gali, J., M. Gertler, and D. Lopez-Salido}\ (forthcoming):

\textquotedblleft Markups, Gaps, and the Welfare Costs of Business

Fluctuations,\textquotedblright\ \emph{Review of Economics and Statistics}.



\textsc{Hall, R.} (1997): \textquotedblleft Macroeconomic Fluctuations and

the Allocation of Time,\textquotedblright\ \textit{Journal} \textit{of Labor

Economics,} 15(Part 2), S223--S250.



\textsc{Hornstein, A., and E. Prescott} (1993): \textquotedblleft The Firm

and the Plant in General Equilibrium Theory,\textquotedblright\ in \textit{%

General Equilibrium, Growth, and Trade, Vol. 2: The Legacy of Lionel McKenzie%

}, ed. by R. Becker et al. San Diego: Academic Press.



\textsc{Ingram, B., N. Kocherlakota, and N. Savin}\ (1997):

\textquotedblleft Using Theory for Measurement: An Analysis of the Cyclical

Behavior of Home Production,\textquotedblright\ \textit{Journal of Monetary

Economics}, 40, 435--456.



\textsc{Kehoe, T., and E. Prescott} (2002): \textquotedblleft Great

Depressions of the 20th Century,\textquotedblright\ \textit{Review of

Economic Dynamics}, 5, 1--18.



\textsc{Kiyotaki, N., and J. Moore} (1997): \textquotedblleft Credit

Cycles,\textquotedblright\ \textit{Journal of Political Economy}, 105,

211--248.



\textsc{Kydland, F., and E. Prescott} (1982): "Time to Build and Aggregate

Fluctuations,\textquotedblright\ \textit{Econometrica}, 50, 1345--1370.



%TCIMACRO{\TeXButton{rule}{\rule[-0.02in]{0.81in}{0.2pt}} }%

%BeginExpansion

\rule[-0.02in]{0.81in}{0.2pt}

%EndExpansion

(1988): \textquotedblleft The Workweek of Capital and Its Cyclical

Implications,\textquotedblright\ \textit{Journal of Monetary Economics}, 21,

343--360.



\textsc{Lagos, R.} (2006): \textquotedblleft A Model of

TFP,\textquotedblright\ \emph{Review of Economic Studies}, 73, 983--1007.



\textsc{McGrattan, E.} (1991): \textquotedblleft The Macroeconomic Effects

of Distortionary Taxation,\textquotedblright\ Institute for Empirical

Macroeconomics Discussion Paper 37, Federal Reserve Bank of Minneapolis.



\textsc{Mulligan, C.}\ (2002a): \textquotedblleft A Century of Labor-Leisure

Distortions,\textquotedblright\ Working Paper 8774, National Bureau of

Economic Research.



%TCIMACRO{\TeXButton{rule}{\rule[-0.02in]{0.81in}{0.2pt}}}%

%BeginExpansion

\rule[-0.02in]{0.81in}{0.2pt}%

%EndExpansion

\textbf{\ }(2002b): \textquotedblleft A Dual Method of Empirically

Evaluating Dynamic Competitive Equilibrium Models with Market Distortions,

Applied to the Great Depression and World War II,\textquotedblright\ Working

Paper 8775, National Bureau of Economic Research.



\textsc{Parkin, M.} (1988): \textquotedblleft A Method for Determining

Whether Parameters in Aggregative Models Are Structural,\textquotedblright\ 

\textit{Carnegie-Rochester Conference Series on Public Policy}, 29, 215--252.



\textsc{Prescott, E.} (1986): \textquotedblleft Theory Ahead of Business

Cycle Measurement,\textquotedblright\ \textit{Federal} \textit{Reserve Bank

of Minneapolis Quarterly Review}, 10(Fall), 9--22.



%TCIMACRO{\TeXButton{rule}{\rule[-0.02in]{0.81in}{0.2pt}} }%

%BeginExpansion

\rule[-0.02in]{0.81in}{0.2pt}

%EndExpansion

(1999): \textquotedblleft Some Observations on the Great

Depression,\textquotedblright\ \textit{Federal Reserve} \textit{Bank of

Minneapolis Quarterly Review}, 23(Winter), 25--31.



\textsc{Restuccia, D., and R. Rogerson} (2003): \textquotedblleft Policy

Distortions and Aggregate Productivity with Heterogeneous

Plants,\textquotedblright\ Manuscript, Arizona State University.



\textsc{Rotemberg, J., and M. Woodford} (1991): \textquotedblleft Markups

and the Business Cycle,\textquotedblright\ in \textit{NBER Macroeconomics

Annual 1991,} ed. by O. Blanchard and S. Fischer. Cambridge, Mass.: MIT\

Press.



%TCIMACRO{\TeXButton{rule}{\rule[-0.02in]{0.81in}{0.2pt}} }%

%BeginExpansion

\rule[-0.02in]{0.81in}{0.2pt}

%EndExpansion

(1999): \textquotedblleft The Cyclical Behavior of Prices and

Costs,\textquotedblright\ in \textit{Handbook of Macroeconomics,} ed. by J.

Taylor and M. Woodford.\textit{\ }Amsterdam: North-Holland.



\textsc{Schmitz, J.}\ (2005): \textquotedblleft What Determines Labor

Productivity? Lessons From the Dramatic Recovery of the U.S. and Canadian

Iron-Ore Industries Following Their Early 1980s Crisis,\textquotedblright\ 

\emph{Journal of Political Economy,} 113, 582--625.



\newpage \begingroup\setlength{\parindent}{0pt}\setlength{\parskip}{0ex}%

\renewcommand{\enotesize}{\normalsize}\theendnotes\endgroup



\end{document}

